\documentclass[11pt]{article}
\usepackage[margin=1in]{geometry}
\usepackage{amsmath,amssymb,amsthm,mathtools}
\usepackage[T1]{fontenc}
\usepackage{lmodern}
\usepackage{microtype}
\usepackage[colorlinks=true,linkcolor=blue,urlcolor=blue,citecolor=blue]{hyperref}
\setlength{\emergencystretch}{3em}
\newtheorem{theorem}{Theorem}[section]
\newtheorem{lemma}[theorem]{Lemma}
\newtheorem{definition}[theorem]{Definition}
\newcommand{\R}{\mathbb R}
\newcommand{\N}{\mathbb N}
\newcommand{\Z}{\mathbb Z}
\newcommand{\C}{\mathbb C}
\newcommand{\code}[1]{{\small\nolinkurl{#1}}}
\title{Polynomial stability and computability\\
\large Proof explanations, formal interfaces, and verified scope}
\author{Companion to the Lean development of Jun Liu and Ivan O. Shevchenko}
\date{31 October 2026}
\begin{document}
\maketitle

\section{Scope and preservation of the published proof}

This companion describes new work in the \code{Stability} library. The
reference points are the published Lean commit \code{81276ab} and the
revised manuscript \cite{paper} at Overleaf commit \code{6f5bced}, in
\texttt{new Main Document}. The published \code{HGAS} library, its public
theorem statements, and its Lean/Mathlib 4.30.0 environment are preserved.
The two earlier companions continue to document their respective releases.

The new library is not a claim that every revised intermediate statement
has already been formalized. The bounded-degree polynomial GAS and LAS
corollaries and uniform dimension-and-odd-degree-input completeness theorem below
are proved. The Fourier estimate, recurrence lemma, restart-budget lemma,
and cubic-relaxation estimates are completed supporting results, with the limits of
their correspondence stated explicitly. The broader source-level
trajectory/test equivalence and the Section~9 compact-base trajectory
extraction are now proved as well. The canonical-graph size and polynomial
Fourier bounds are explicit. The full halting-to-HLS reduction and its
cubic/cofinal hardness consequences are now proved. The printed gain and
reward have their exact two-sided invariant-average bounds: the positive
branch constructs a probability on the actual compact graph, and the
negative branch uses ordinary-time extraction. The printed-gain compiler
is computable and gives the fixed-pair and cofinal GAS completeness
theorems. The existing alternative-gain compiler remains unchanged.
\code{FORMALIZATION_LEDGER.md} records the current status.

The simulation estimates follow the written proof in \cite{korda}; the
author's earlier Lean implementation is linked in the README. No code is
copied or ported from it.\footnote{Parts of that implementation were inspected
during an earlier feasibility review; no clean-room claim is made.}

\section{Local attraction becomes global under homogeneity}

The formal definition of a forward trajectory includes its initial value
and the differential equation on the entire nonnegative time ray. Merely
quantifying over such trajectories can be vacuous when a solution blows up
in finite time. We therefore include local existence explicitly.

\begin{definition}
An equilibrium $x_*$ is locally asymptotically stable if it is Lyapunov
stable with local existence of global forward trajectories, and there is
$r>0$ such that every initial state in $B_r(x_*)$ admits a global forward
trajectory and every such trajectory tends to $x_*$.
\end{definition}

Here stability with existence means that for every $\varepsilon>0$ there
is $\delta>0$ such that every initial state in $B_\delta(x_*)$ admits a
global forward trajectory, and every such trajectory stays in
$B_\varepsilon(x_*)$. This is the nonvacuous convention already used by
the published HLS development. For locally unique polynomial ODEs it agrees
with the usual initial-value formulation. The theorem below is proved for
the stated trajectory predicates and does not assume a uniqueness theorem.

\begin{theorem}\label{thm:lasgas}
Let $E$ be a real normed vector space, let $k\in\N$, and suppose
\[
 F(cx)=c^{k+1}F(x)\qquad(c\in\R,\ x\in E).
\]
Then the origin is locally asymptotically stable if and only if it is
globally asymptotically stable.
\end{theorem}

\begin{proof}
Global stability and attraction imply their local counterparts, so consider
the other direction. The published homogeneous stability lemma gives a
global forward trajectory from every initial state: shrink the state into
a neighborhood where existence is known, then scale its trajectory back.
The same lemma gives the global linear trajectory bound equivalent to
Lyapunov stability with existence. For the present theorem we need its
existence conclusion; the local stability hypothesis already gives stability.

It remains to prove attraction from an arbitrary initial state $x_0$.
Let $r>0$ be a local attraction radius, and put
\[
 c=\frac{r}{2(\lVert x_0\rVert+1)}>0.
\]
Then $\lVert cx_0\rVert<r$. If $\gamma$ is any forward trajectory from
$x_0$, homogeneity gives a forward trajectory
\[
 \widetilde\gamma(s)=c\gamma(c^ks)
\]
from $cx_0$. Given $\varepsilon>0$, local attraction gives $T\ge0$ with
$\lVert\widetilde\gamma(s)\rVert<c\varepsilon$ for all $s\ge T$.
For $t\ge Tc^k$, set $s=t/c^k$. Then
\[
 c\lVert\gamma(t)\rVert=\lVert\widetilde\gamma(s)\rVert
 <c\varepsilon.
\]
Dividing by $c>0$ proves attraction. The choice of $c$ also handles
$x_0=0$, so no separate zero-state argument is needed.
\end{proof}

The definitions and proof are in
\code{Stability/Dynamics/LocalAsymptoticStability.lean}. The main
declaration is
\code{locallyAsymptoticallyStable_iff_globallyAsymptoticallyStable}.
It reuses the published trajectory-scaling theorem and homogeneous
linear-bound characterization \cite[Lemma 1.5]{korda}.

\section{General polynomial GAS and LAS}

The input is a finite sparse rational code for a polynomial vector field.
Our new languages \code{PGASQ n d} and \code{PLASQ n d} require each decoded
coordinate polynomial to have total degree at most $d$, require the origin
to be an equilibrium, and require GAS or LAS, respectively. The degree
bound is imposed \emph{after decoding}: two entries of arbitrarily high
degree may cancel. A bound on the largest syntactic exponent would describe
a different input condition.

\begin{theorem}
Let $N_*$ be the fixed dimension threshold of the published cubic HGAS
completeness theorem. For every fixed $n\ge N_*$ and $d\ge3$, both
\code{PGASQ n d} and \code{PLASQ n d} are $\Sigma^0_1$-hard. Each is
undecidable, and its complement is not recursively enumerable.
\end{theorem}

\begin{proof}
Homogeneous cubic GAS is $\Sigma^0_1$-complete in dimension $n$ by the
published theorem. A reduction must nevertheless handle every raw input,
including codes that are not homogeneous cubic fields. Define
\[
 C(F)=\begin{cases}
 F,&\text{if every decoded coordinate of $F$ is homogeneous of degree 3},\\
 0,&\text{otherwise}.
 \end{cases}
\]
The homogeneity test is primitive recursive, so this is a total primitive-
recursive map. For a valid input, its decoded degree is at most three and
the GAS assertion is unchanged. For an invalid input, $C(F)$ is the zero
field. In positive dimension this field is not GAS: a nonzero constant
trajectory does not tend to the origin. The existing threshold $N_*$ is
positive, so this case applies throughout the theorem.

For LAS, Theorem~\ref{thm:lasgas} identifies LAS and GAS on every valid
homogeneous cubic input. It also excludes LAS of the zero fallback, which
is homogeneous of every positive degree. Thus the same map $C$ is a
many-one reduction to both target languages. No parity restriction on $d$
is needed, because $d$ is an upper bound rather than a specified homogeneous
degree. Undecidability and failure of recursive enumerability of the
complements follow by pulling a hypothetical decision or complement
semidecision procedure back along the reduction.
\end{proof}

These are hardness results, not completeness results: this argument does
not show that general polynomial GAS or LAS is recursively enumerable.
The two modules are
\begin{quote}
\code{Stability/Computability/PolynomialGAS.lean}\\
\code{Stability/Computability/PolynomialLAS.lean}.
\end{quote}
Their correctness declarations include both branches of the raw-code
reduction:
\begin{quote}
\code{cubicToPolynomialCode_correct}\\
\code{cubicToPolynomialCode_correct_las}.
\end{quote}
The endpoint names begin with \code{pgasQ_} and \code{plasQ_}.

\section{A single checker when the degree is input data}

For fixed dimension $n$, the language \code{VariableOddHGASQ n} takes a pair
$(d,F)$ as input and requires $d$ to be odd and $F$ to belong to
\code{HGASQ n d}. One must not infer its recursive enumerability merely
from the existence of a semidecision procedure for each fixed $d$.

Instead, \code{Stability/Computability/VariableHGASDegree.lean} constructs
one primitive-recursive Boolean relation on $(d,F,c)$, where $c\in\N$
encodes a rational Lyapunov candidate. It checks parity and homogeneity,
then runs the published exact-grid positivity tests on the candidate and
its derivative numerator. The candidate type and its decoder depend on
$n$, but not on $d$. The published soundness theorem and rational homogeneous
converse show that an accepted witness exists exactly for inputs in the
semantic language. Searching the natural witnesses therefore proves
\code{variableOddHGASQ_re} uniformly in the degree input.

The computable map $F\mapsto(3,F)$ embeds the fixed cubic language. Thus
\code{variableOddHGASQ_sigmaOneComplete} proves completeness for every
fixed $n\ge N_*$. This is a degree-input result with an explicit oddness
condition.

\subsection{A uniform checker with dimension as input}

\code{VariableDimensionHGAS.lean} now supplies the further uniformity.
An input is $(n,d,F)$, where $F$ is a list of sparse polynomials and each
exponent vector is itself a list of natural numbers. The language
\code{DimensionDegreeHGASQ} requires exactly $n$ field coordinates,
exactly $n$ entries in each exponent vector, odd $d$, and the existing
semantic predicate \code{HGASQ n d} for the decoded field. Padding makes
decoding total, but malformed shapes are rejected before this can affect
accepted semantics. Coefficients of duplicate monomials are added exactly,
so cancelling terms do not create false degree restrictions.

A witness has a dimension-independent type: a natural exponent, one raw
polynomial, and two natural grid resolutions. The raw arithmetic modules
implement derivatives, the rational Lyapunov derivative numerator, and
the grid tests directly with list operations. Cartesian grid words are
generated by recursion on the input dimension. These operations form one
jointly primitive-recursive Boolean checker, proved to commute exactly
with the old fixed-dimensional checker on well-formed codes. An unsuccessful
natural-number witness decode produces zero grid resolutions, which are
rejected even in dimension zero.

Write $\mathcal H$ for \code{DimensionDegreeHGASQ} and $C$ for
\code{RawHGASChecker.check}. The published soundness and converse theorems give
\[
 (n,d,F)\in\mathcal H
 \quad\Longleftrightarrow\quad
 \exists c\in\N:\ C((n,d,F),c)=\mathrm{true}.
\]
This proves uniform recursive enumerability. The computable encoding of
one fixed cubic slice gives $\Sigma^0_1$-completeness, undecidability, and
failure of recursive enumerability of the complement. This language
deliberately rejects even degree tags; no classification of those tags
is inferred from the odd-degree converse theorem. The regression tests
include dimension zero, malformed shapes, failed-resolution witnesses,
and an accepted certificate for $-x^3$ displayed with cancelling quadratic
terms.

These results give partial coverage of Corollary \code{cor:variable} in
the manuscript. Separate dimension-input languages with fixed degree
three or fixed odd degree are not formalized. No literal identification
with the paper's unrestricted all-degree language is asserted: the joint
Lean language admits dimension zero and explicitly requires odd degree.

\section{Shifted theta exponents and weighted summability}

Step 1 of the revised manuscript's Fourier bound uses the shifted exponent
\[
 \sigma(n)=(n+1)^2-\tfrac14.
\]
It is negative at $n=-1$. Treating every exponent as nonnegative would lose
that term. A Gaussian estimate with arbitrary linear and constant parts
handles it directly, without a special exception.

\begin{lemma}\label{lem:gaussian}
Suppose $a>0$, $\lambda>0$, $\rho\ge1$, $A\ge0$, and
$\lambda\rho^A<1$. For arbitrary $b,c,B,C\in\R$ and $d\in\N$, the series
\[
 \sum_{n\in\Z}(1+|n|^d)
 \lambda^{a n^2+bn+c}
 \rho^{A|a n^2+bn+c|+B|n|+C}
\]
converges.
\end{lemma}

\begin{proof}
Take logarithms of the two exponential factors. Since $\log\rho\ge0$,
\[
 |an^2+bn+c|\le an^2+|b||n|+|c|
\]
gives an upper bound of the form
\[
 -\alpha n^2+\beta|n|+\gamma,
 \qquad \alpha=-a(\log\lambda+A\log\rho)>0.
\]
The signed linear term $(\log\lambda)bn$ is at most
$|b\log\lambda||n|$. Thus the summand is bounded by a fixed polynomial
factor times a Gaussian with an arbitrary linear shift. The existing
Gaussian summability theorem proves convergence of this majorant.
Every inequality holds for every integer $n$, including those for which
$an^2+bn+c<0$.
\end{proof}

For $\lambda>0$, $\rho\ge1$, and $D\in\R$, take $a=1$, $b=2$,
$c=3/4$, $A=4$, $B=4D$, $C=2D$, and $d=1$. The lemma and
$4(2|n|+1)\le8(1+|n|)$ prove convergence of
\[
 \sum_{n\in\Z}4(2|n|+1)\lambda^{(n+1)^2-1/4}
 \rho^{4|(n+1)^2-1/4|+2(2|n|+1)D}
\]
whenever $\lambda\rho^4<1$. This is an all-index majorant corresponding to
the revised coefficient and frequency bounds; the printed proof instead
separates the finitely many negative exponents. A comparison theorem permits any complex coefficients
and frequencies satisfying the corresponding bounds.

These results are in
\code{Stability/Analysis/ShiftedThetaBounds.lean}. The actual graph
correspondence additionally requires collecting equal frequencies, finite
products, and reciprocal expansions. The later graph and Fourier modules
prove those identities and give the explicit coordinate budgets; the
scalar majorant alone is not a substitute for that correspondence.

\section{Why the intended torus trajectories are recurrent}

\begin{lemma}
Let $G$ be a compact first-countable topological abelian group and let
$f:\R\to G$ be an additive homomorphism. For every $x\in G$, the curve
$t\mapsto x+f(t)$ returns to its initial point along times tending to
$+\infty$. So does every continuous observation of this curve.
\end{lemma}

\begin{proof}
Compactness gives a strictly increasing sequence $\phi:\N\to\N$ such
that $f(\phi(n))$ converges to some $y\in G$. Set
$\tau_n=\phi(2n)-\phi(n)$. Strict increase implies $\tau_n\ge n$, so
$\tau_n\to+\infty$. Additivity and continuity of subtraction give
\[
 f(\tau_n)=f(\phi(2n))-f(\phi(n))\longrightarrow y-y=0.
\]
Hence $x+f(\tau_n)\to x$. Applying a continuous observation preserves
this convergence. No nonresonance assumption is needed.
\end{proof}

The generic result and its finite-torus specialization are in
\begin{quote}
\code{Stability/Dynamics/CompactTranslationRecurrence.lean}.
\end{quote}
This proves
the recurrence step in the revised intended-trajectory argument. It does
not assert the full source ODE, validity residuals, or test averages for
that graph; those must be connected separately.

\section{Forward-only ingredients for the broader selector}

The revised selector assumes forward recurrence and bounded forward
averages of a real power $\mu\ge1$. Two points require care when reusing
the published formalization: a negative clock must not introduce a
reverse-recurrence assumption, and a bound at exponent one cannot be
replaced by a bound at exponent two.

\subsection{Changing the sign of a clock without reversing time}

The unary modular vector field has rational coefficients. Complex
conjugation therefore takes a trajectory with clock $i\Omega$ to one
with clock $-i\Omega$. It preserves every coordinate magnitude and the
same forward return-time sequence. If $\Omega<0$, conjugate first; then
rescale time by the positive factor $\pi/|\Omega|$. The resulting clock
is the standard $i\pi$, with forward recurrence intact.

Conjugation commutes with each finite-horizon Ces\`aro mean. It also
changes the first-harmonic factor $e^{-i\alpha}$ to $e^{i\alpha}$, so
the corresponding phase is $-\alpha$ and real target moments are
unchanged. Positive time rescaling preserves the full forward limits and
transports norm-power averages by rescaling their horizons.

These are exact statements in
\code{Stability/Selection/UnaryClockConjugation.lean}. The linked
Jacobi module \code{LinkedClockConjugation.lean} conjugates all residuals
and both clocks. An important qualification is that a real elliptic clock
is unchanged under conjugation, while a purely imaginary modular clock
changes sign. We do not identify this transformation with simultaneous
negation of all physical Jacobi frequencies.

\subsection{The critical linear cusp estimate}

In a cusp, Haar measure has a positive height marginal bounded below by
a multiple of $dy/y^2$. The first height moment already diverges:
\[
 \int_A^\infty y\frac{dy}{y^2}=\int_A^\infty\frac{dy}{y}=\infty.
\]
The module \code{Selection/LinearCuspDivergence.lean} proves this with the
published cusp measure and arbitrary finite cutoffs. Taking the square root
of the published nonnegative square-theta observable gives a continuous
linear observable on the modular frame quotient. Its cusp lower bound
implies infinite Haar integral.

The existing unconditional pointwise horocycle theorem applies in either
direction of nonzero speed. On a nonperiodic orbit, continuous compactly
supported minorants with arbitrarily large Haar integrals force the linear
observable's forward averages to diverge. Along the reconstructed modular
trajectory, that observable is exactly the norm of the $C$ coordinate.
For every $u\ge0$ and real $\mu\ge1$, the elementary inequality
$u\le1+u^\mu$ turns an upper bound on the $\mu$-power averages into an
upper bound on linear averages. Hence such a trajectory cannot have a
nonperiodic frame orbit.

This exclusion is formalized in
\code{Stability/Selection/LinearThetaCuspExclusion.lean}. The endpoint
assumes the normalized unary ODE, nonvanishing $B$, the lambda-coordinate
identity, and an actual reconstructed frame path. It proves the
nonperiodic-orbit exclusion. The next modules supply the lift and frame
from the recurrent trajectory itself.

\subsection{Selecting the unary trajectory from forward means}

Three modules combine the preceding ingredients:
\begin{quote}
\code{OrdinaryTimeGenericity.lean}\\
\code{OrdinaryTimeClosedFrame.lean}\\
\code{OrdinaryTimeUnarySelection.lean}.
\end{quote}
The strongest endpoint starts with an all-real-time linked trajectory
whose modular clock is $i\pi$, one forward recurrence sequence of the
literal linked polynomial state, all ten linked residuals, and the limits
\[
 \langle Y\rangle=1,\qquad \langle p_4\rangle=1,\qquad
 \langle Y(t)\eta e^{-i\pi t}\rangle=\frac1{128},\qquad |\eta|=1.
\]
Here brackets mean the full limit of $T^{-1}\int_0^T$ as real $T\to\infty$.
The remaining hypothesis is an eventual upper bound on the corresponding
averages of the full linked-state norm to a real power $\mu\ge1$.

The nonzero harmonic mean excludes the three degenerate divisors.
Indeed, recurrence and the unary ODE force $Y$ to be constant on either
of the divisors $A=0$ and $C=0$. On $B=0$, the residual $Y^4=B$ gives
$Y=0$. A constant has zero mean against the nonzero-frequency character,
contrary to the stipulated limit. No auxiliary $p_2$ moment is needed.

On the generic locus, the published reconstruction theorem produces a
differentiable lambda lift and either a horizontal path or an exact frame
path. The full-state norm controls the unary $C$ coordinate, so the linear
cusp argument excludes the nonperiodic frame case. A closed chart follows.
Its cusp-amplitude classification and the two raw means $\langle Y\rangle$
and $\langle p_4\rangle$ select the infinity cusp and unit speed. Finally,
the first harmonic is
\[
 2\eta e^{-\pi y}e^{i\pi x}=\frac1{128}.
\]
Taking absolute values fixes the gauge height $y=y_*$ with
$e^{-\pi y_*}=1/256$; the same identity then gives
$\eta e^{i\pi x}=1$.

Thus the entire unary state is the canonical theta state at
$x+t+iy_*$, and $p_4(t)=\theta_4(x+t+iy_*)^2$. The frame and genericity
conditions are conclusions, not assumptions. This linked result still
uses the $p_4$ mean and all linked residuals; it is not a claim to have
formalized the standalone unary proposition with only two means.

The module \code{OrdinaryTimePhysicalClock.lean} handles a general clock
$i\Omega$, $\Omega\ne0$, with recurrence at time zero. For $\Omega>0$,
the substitution $t=(\pi/\Omega)s$ preserves forward horizons and gives
the normalized theorem. For $\Omega<0$, first conjugate, apply the
positive-clock theorem, and use the series identity
\[
 \overline{\theta_j(\tau)}=\theta_j(-\overline\tau)
\]
to return to the original state. The conclusion has the signed speed
$\Omega/\pi$ and the same literal phase equation. No step samples the
original curve at negative times. The elliptic clock is conjugated and
positively rescaled according to its own value; a real elliptic clock is
not negated by conjugation.

\subsection{The gauged Jacobi output}

The linked residuals recover $p_3$ and the Eisenstein coordinate from the
canonical unary state. They recover $p_2$ up to sign; translating the
modulus by two absorbs that sign while preserving its height and the
other roots. This supplies a canonical initial root state. The published
theta parametrization and reciprocal-chart continuation theorem then
produce one globally regular affine Bloch path. The linked forward return
sequence projects to the Jacobi coordinates, so the published recurrence
criterion forces the Bloch parameter to be real. These steps are combined
in \code{OrdinaryTimeBlochClassification.lean} without assuming a Bloch
representation at the outset.

For physical clocks $(\Omega+N)/2$ and $i\Omega$, write
$\rho=N/\Omega$ and assume $\Omega\ne0$ and $\rho$ irrational. The gauge
ODE
\[
 g'=\frac{i(\Omega+2N)}4\,g
\]
determines the actual gauge as its initial value times an exponential;
no normalization of $g(0)$ is imposed. The gauged output $Z=gT_1$ has the
absolutely summable quadratic Fourier representation of that Bloch path.

Let $a_n=R^{n^2}S^n$, with $R,S\ne0$, be a target coefficient family, and
let $c_+,c_-\ne0$ be the initial character multipliers. The module
\code{OrdinaryTimeJacobiMoments.lean} uses the three further forward means
\[
 \langle Z\rangle=1,\qquad
 \langle c_+e^{-i(\Omega+N)t}Z\rangle=c_+a_1,\qquad
 \langle c_-e^{-i(\Omega-N)t}Z\rangle=c_-a_{-1}.
\]
The published three-moment rigidity theorem identifies the entire actual
output with the quadratic trace of the target coefficients. The proof
derives the representation, its summability, and the gauge branch from
the trajectory data. It does not assume a selected invariant measure or
backward recurrence.

\subsection{The literal source coordinates and arithmetic detector}

The modules from \code{OrdinaryTimeDetectorProjection.lean} through
\code{OrdinaryTimeDetectorAverage.lean} apply this classification to the
actual rational polynomial source field. Pointwise validity identifies
its linked residuals, reciprocal variables, and unit clocks. Two extra
residuals in the compiler follow algebraically from the eight displayed
relations; they do not impose a stronger source condition.

The controller equations and forward recurrence force constant clock
parameters. For distinct primes $p_j$, these parameters give signed
frequencies $\pm\sqrt{p_j}$. Integer independence survives both those
signs and the common scaling factor. The nonzero-column condition on the
integer input then gives the irrational ratios needed above. Thus these
ratios, the phase trajectory, and the target theta series are conclusions,
not additional assumptions on the source trajectory.

Let $\Phi$ denote the continuous arithmetic Fourier detector on the
phase torus. The actual polynomial detector $D(x(t))$ satisfies
\[
 D(x(t))=c\,\Phi(\varphi_0+t\omega),\qquad |c|=1.
\]
Here the reciprocal graph equations supply the Poisson factors, and the
actual right-hand-side phase ODE supplies the constant $c$. Absolute
Fourier summability and integer independence give the full forward limit
\[
 \lim_{H\to\infty}\frac1H\int_0^H D(x(t))\,dt
 =c\int\Phi.
\]
The published arithmetic constant-term theorem says that $\int\Phi=0$
exactly when the integer system has no solution. Since $c\ne0$, both
real scalar detector means vanish exactly in that case.

This theorem concerns a \emph{given qualifying trajectory}. Its hypotheses
are the literal ODE for all real times, pointwise validity, forward
recurrence at zero, the specified scalar moment means, and an eventual
bound on the full Euclidean $\mu$-moment for some real $\mu\ge1$.
No backward recurrence or selected invariant measure is assumed.
Pointwise validity is not inferred from zero residual averages.
The compact-base extraction described below supplies these hypotheses
from the measure assumptions in Section~9.

\subsection{Constructing a trajectory for the converse}

The modules from \code{OrdinaryTimeIntendedOrbit.lean} through
\code{OrdinaryTimePrimeExistence.lean} now supply the other direction.
Start with the closure of the published bounded canonical orbit. Its
compact dynamics admit an invariant probability, from which the proof
constructs a forward recurrent point. The positive physical clock has a
uniform lower speed bound; its inverse therefore defines a complete
ordinary-time source trajectory and transports the forward return times.
Continuity extends the validity identities and uniform state bound from
the dense canonical orbit to the closure.

For each actual scalar drift coordinate, the published Fourier
construction provides a bounded primitive on that closure. Before the
clock change its derivative is the clock speed times the scalar drift.
Differentiation after the inverse clock change cancels that speed. Thus
the ordinary-time integral is the difference of two bounded primitive
values. Dividing by the full real horizon $H\to\infty$ proves a zero
mean. No desired test mean is assumed. The uniform state bound also
controls every nonnegative real moment at every positive horizon.

Combining this construction with the exclusion direction gives
\code{exists_prime_source_trajectory_iff_not_solvable}: for distinct
prime clocks, nonzero linear columns, and every fixed real $\mu\ge1$, a
complete valid forward-recurrent trajectory satisfying the growth and
all scalar test conditions exists if and only if the integer system has
no solution. The corrected compiler centerings equal the printed
centerings under the nonzero-column hypothesis; that equality is proved
for the literal scalar codes before combining the two directions.

This converse constructs an orbit-closure trajectory rather than the
manuscript's literal torus graph $\Gamma$. It proves the source-level
existence statement, but does not by itself identify all of that graph's
weighted Fourier norms.

\subsection{Extracting the trajectory from a compact-base measure}

\code{OrdinaryTimeDetectorExtraction.lean} starts with an ergodic invariant
probability for the actual compact base flow. It assumes full mass in the
finite chart and zero integrals of every compact scalar drift, with unit
gain. It then constructs a complete ordinary-time source trajectory,
forward recurrent at zero, satisfying pointwise validity and all full
real-horizon scalar mean conditions. The growth exponent is the actual
even denominator-clearing exponent $\mu$, and the bound holds at every
horizon $H\ge1$. The proof also checks $\mu\ge1$.

The construction first restricts to the finite chart and selects a point
that is both recurrent and generic for the needed observables. A positive
clock change converts its compact orbit to an ordinary-time orbit. The
clock-speed averages in both time directions ensure that this new orbit
exists for all real times; no backward recurrence is assumed. A
nonnegative validity-energy coordinate has zero integral, so continuity
and invariance give its pointwise vanishing along the selected orbit.
The other zero means survive the clock change by the corresponding
ratio-of-integrals argument. An average estimate for the clock weight
gives the Euclidean growth bound.

Thus the recurrence, validity, and ordinary-time averages are conclusions,
not extra selection hypotheses. Combining this extraction with the source
equivalence shows that a solvable arithmetic instance admits no such
ergodic finite-chart probability. \code{OrdinaryTimeConditionalBase.lean}
obtains precisely this base measure from an ergodic accumulator measure
with positive finite-interior mass: restriction and the equivariant base
projection preserve ergodicity, the factor has full finite-chart mass,
and a nonzero drift gain cancels from the zero-mean equations.

\code{OrdinaryTimeAccumulatorGap.lean} combines the resulting solvable-
instance obstruction with the published nonsolvable-instance finite graph.
This proves the complete compact sign gap for the actual effective compiler,
and its homogeneous wrapper is GAS exactly when the arithmetic instance
is solvable. \code{OrdinaryTimeUniversalFamily.lean} supplies nonzero
columns by the existing computable augmentation and gives fixed-pair and
cofinal $\Sigma^0_1$-completeness on this proof route. The output code is
definitionally the existing finite compiler; the selected measure is not
an input to code generation.

This first construction retains the published effective gain
$1/[8(T+1)(R+1)]$ and reward margin. The following modules separately
implement Section~9's least-integer gain, chosen by $m^2w_*\ge4K_e^2$,
and margin $w_*/4$. They prove the full torus-graph norm package and the
two-sided invariant-average bounds for those literal constants. The two
code families need not be equal; each has its own checked sign argument.

\subsection{An actual weighted Fourier coordinate}

\code{OrdinaryTimeWeightedFourier.lean} identifies the Haar Fourier
coefficients of an absolutely summable character expansion by collecting
all terms with the same frequency. Its weighted norm is bounded by the
sum of the weighted expansion coefficients, even when frequencies collide.
This avoids assuming that a formal expansion is already the Fourier
series with distinct indices.

\code{OrdinaryTimeThetaGraph.lean} applies this result to the half-period
theta coordinate $T_1$ of the literal torus graph. In fundamental clock
coordinates, its $n$-th term has coefficient $4\lambda^{n^2}$, with
$\lambda=1/256$, and frequency length
\[
 |4n^2-1|+|4n-2|\sum_\ell |B_{\ell j}|.
\]
The function agrees with the shifted Jacobi value at every real lift of
the torus phases, not just along one orbit. The corrected Gaussian bound
proves summability of its actual weighted Haar coefficients whenever
$\lambda\rho^4<1$ and $\rho\ge1$. The module also computes the exact
right-hand-side character norm $\rho^{4\sum_\ell|b_\ell|}$.

\code{OrdinaryTimeFourierInverse.lean} proves a reusable reciprocal
estimate. If an absolutely summable character expansion $h$ has weighted
coefficient sum at most $\delta<1$, expand
$(1-h)^{-1}=\sum_{r\ge0}h^r$. Each power is indexed by finite tuples of
the original terms. Submultiplicativity of the weight bounds its sum by
$\delta^r$. Collecting equal frequencies only decreases the resulting
upper bound, so the actual weighted Haar norm is at most $(1-\delta)^{-1}$.
This proves a statement about the reciprocal function, not just a formal
series with an assumed norm.

\code{OrdinaryTimePoissonGraph.lean} applies this argument to both signs
of each Poisson reciprocal. The norm is at most $4$ when the corresponding
weighted geometric ratio is at most $3/4$.
For the theta reciprocal, remove the leading gauge character and factor
$T_1$ as $4$ times a character times a centered theta function. Under
\[
 \frac1{256}\rho^{4+4\sum_\ell|B_{\ell j}|}\le\frac18,
\]
the nonconstant centered tail has weighted norm at most $2/7$.
Its reciprocal therefore has norm at most $7/5$. Restoring the gauge
character gives the literal coordinate $U=1/T_1$ and the bound
\[
 \|U\|_\rho\le\frac7{20}\rho^{1+2\sum_\ell|B_{\ell j}|}.
\]
\code{OrdinaryTimeThetaReciprocal.lean} proves nonvanishing and the
identity $T_1U=1$ at every phase. \code{OrdinaryTimeJacobiGraph.lean}
also identifies all four Jacobi coordinates with explicit torus series
and proves their actual weighted summability. The derivative, unary theta,
and Eisenstein modules provide the remaining complex coordinates.

\subsection{Assembling the graph and composing with polynomials}

\code{OrdinaryTimeSourceGraph.lean} assembles these coordinates into the
literal source-state type and its real coordinate encoding. It proves
continuity of the resulting map $\Gamma$. The Fourier algebra modules
work with actual Haar coefficients: products satisfy the weighted product
bound, and real or imaginary parts do not increase the weighted norm.
The latter fact follows from reflection of frequencies under complex
conjugation. For example,
\[
 \widehat{\operatorname{Re}f}(a)
 =\frac{\widehat f(a)+\overline{\widehat f(-a)}}2.
\]
Thus realification does not introduce an unnecessary factor of two.

\code{OrdinaryTimeSourceFourier.lean} combines the coordinate estimates.
Every actual coordinate of $\Gamma$ has summable weighted Fourier
coefficients under the two printed matrix-norm conditions on $\rho$.
Changing the arithmetic right-hand side $b$ changes only its unit-modulus
oscillator. In particular, the squared Euclidean norm of $\Gamma$ is
exactly independent of $b$, not just bounded uniformly in $b$.

The polynomial composition lemma has a useful general form. If each
coordinate function has weighted norm at most $C\ge1$, then
\[
 \|P\circ\Gamma\|_\rho
 \le \left(\sum_e |P_e|\right) C^{\deg P}.
\]
Expand $P$ into monomials, apply the product bound to each monomial, and
use $C\ge1$ to replace its degree by $\deg P$. This proof includes
constant and zero polynomials and does not require distinct frequencies
in the coordinate expansions.

Separating the right-hand-side oscillator gives a natural constant $C$,
uniform in $b$ and $P$, such that
\[
 \|P\circ\Gamma_b\|_\rho
 \le \left(\sum_e |P_e|\right) C^{\deg P}
       \rho^{4\deg P\,\sum_\ell |b_\ell|}.
\]
This statement is proved in \code{OrdinaryTimeSourceFourierBound.lean}.
That first proof uses Archimedean existence. Three numerical modules now
give an explicit choice:
\begin{itemize}
\item \code{OrdinaryTimeThetaBudget.lean}: scalar theta estimates;
\item \code{OrdinaryTimeSourceBudget.lean}: all graph coordinates;
\item \code{OrdinaryTimeSourcePolynomialBudget.lean}: polynomial substitution.
\end{itemize}
Their common coordinate bound is
\[
 C=262144^4+1048576+385+32+1.
\]
This deliberately coarse bound is independent of the arithmetic input.
The scalar estimates bound each Jacobi theta coordinate by $262144$, its
normalized derivative by $1048576$, and the Eisenstein coordinate by
$385$. Polynomial powers, reciprocal coordinates, and the real coordinate
encoding then give the displayed $C$. These estimates retain the signed
shift in the theta exponent, including its negative values.

Fourier reconstruction bounds the absolute value of each coordinate by
its weighted coefficient norm. Applying this at $\rho=1$ gives the
explicit squared Euclidean bound
\[
 \|\Gamma_b(x)\|_2^2\le N C^2,
\]
where $N$ is the source-state dimension. The normalized positive
controller equilibria satisfy the required coordinate bounds because
their two squared coordinates sum to one. Thus neither this size bound
nor the polynomial Fourier bound relies on a non-effective choice of
constant.

\subsection{An explicit primitive bound}

Let the torus have $r$ coordinates with distinct prime clocks $p_i$ and
put $s=2^r-1$. At a radius $\rho=1+1/m$, $m\ge1$, a usable natural
small-divisor multiplier is
\[
 B_{\mathrm{rad}}=
 \left(1+\sum_i p_i\right)^s s!(2m)^s.
\]
The algebraic small-divisor estimate costs a polynomial in the integer
frequency. The elementary binomial estimate
\[
 (1+n)^s\le s!(2m)^s(1+1/m)^n
\]
absorbs that cost into the exponential Fourier weight. This includes
$r=0$, $s=0$, and $n=0$.

For a real polynomial family $(P_j)$ of degrees at most $D$, define
\[
 K_0=1+B_{\mathrm{rad}}\left(\sum_j\sum_e |(P_j)_e|\right)C^D.
\]
The literal Fourier primitive $H$ is continuous and real-valued, and
\[
 \|H(x)\|_2^2\le
 \left(K_0\rho^{4D\sum_\ell|b_\ell|}\right)^2.
\]
The estimate first bounds the sum of the coordinate absolute values,
then uses $\sum_j |H_j|^2\le(\sum_j|H_j|)^2$. No extra dimension factor
is concealed in $K_0$.

Along the physical signed-prime-root flow, its derivative is exactly
$P_j\circ\Gamma_b-\widehat{P_j\circ\Gamma_b}(0)$. The normalization of
the torus frequency by $1/(2\pi)$ cancels the $2\pi$ in differentiating a
character. Conjugate symmetry of the actual Haar coefficients proves
that the primitive is real. The constant coefficient is retained until
its arithmetic vanishing is established; it is not an implicit premise
of the primitive theorem. These results are checked in
\code{OrdinaryTimeRadicalFourierBudget.lean},
\code{OrdinaryTimeRadicalPrimitive.lean}, and
\code{OrdinaryTimeSourcePrimitive.lean}.

\subsection{The rational radius and least gain denominator}

The radius search does not compare arbitrary real numbers. For
$\rho=1+1/m$, its two tests are equivalent to the integer inequalities
\[
 (m+1)^{K_B}\le32m^{K_B},\qquad
 2(m+1)^{K_A}\le3m^{K_A},\qquad m>0.
\]
The rational powers tend to one as $m\to\infty$, so a successful test
exists. The module \code{OrdinaryTimeRadiusChoice.lean} proves that
searching for the least successful integer is a total computable
operation. It also proves that the resulting radius satisfies the
matrix-norm inequalities used above.

For the actual decoded detector drift, the modules
\code{OrdinaryTimeDetectorResources.lean} and
\code{OrdinaryTimeDetectorGain.lean} define the semantic degree $D_h$,
the positive rational $K_0$, and
$K_e=K_0\rho^{4D_h\sum_\ell|b_\ell|}$. The drift family, its degree,
and $K_0$ are unchanged when only the arithmetic right-hand side changes.
Set $M_0=NC^2$ and $w_*=1/(1+M_0)$. The gain module implements the
least positive integer $k$ such that $4K_e^2\le k^2w_*$, using integer
comparisons on a rational code. It proves both leastness and
computability, and obtains
\[
 \kappa=1/k>0,\qquad
 \kappa^2K_e^2\le w_*/4,\qquad \varepsilon=w_*/4>0.
\]
On the source and scaled-primitive graph, write $z=\kappa H(x)$.
The central coordinate satisfies $w^2\ge w_*$, while
$|z|^2\le w_*/4$. Consequently the literal stereographic reward obeys
\[
 \frac{1-|z|^2}{1+|z|^2}+2w^2-1-\varepsilon
 \ge 2w_*-2|z|^2-\varepsilon
 \ge\frac54 w_*.
\]
This is a pointwise bound on the specified graph. Its dynamical use
requires the separate graph-invariance and arithmetic zero-mode results.
The corresponding input algorithm is checked in
\code{OrdinaryTimeDetectorGainCode.lean}. With the matrices and prime tuple
fixed, it computes the budget code primitive-recursively from the integer
right-hand side, then applies the computable least-denominator search.
Equal decoded rational budgets give the same least denominator, so the
proof does not require the intermediate rational codes to be in reduced form.

\subsection{The graph ODE and its arithmetic zero mode}

The explicit graph must solve the source ODE, not just satisfy its
Fourier estimates. The source-flow modules identify every coordinate of
$\Gamma(\varphi+\omega t)$ with the established canonical trajectory at
the same initial phase. This includes the gauge, both reciprocal signs,
and the right-hand-side oscillator. The identity $T_1U=1$ proves the
required nonvanishing chart condition. For the positive prime controllers,
their angular velocity is exactly $\sqrt{p_i}$; its coordinate in
fundamental turns is $\sqrt{p_i}/(2\pi)$. Thus the graph equation uses the
actual polynomial vector field and physical time.

To identify the constant coefficient, assume the arithmetic instance is
unsolvable. The previously checked arithmetic Fourier expansion gives a
bounded primitive for each actual drift coordinate on the zero-phase
orbit, so its full forward time average is zero. Independently, the new
actual-Haar primitive has derivative $f-\widehat f(0)$ and is bounded.
Its derivative also has zero forward mean. Subtracting these two
identities forces $\operatorname{Re}\widehat f(0)=0$; reality of $f$
gives the imaginary part as zero. This proves vanishing of the actual
Haar coefficient, not merely of a coefficient in a formal expansion.
Consequently the primitive equation is uncentered at every initial phase.

The graph-flow endpoint is in \code{OrdinaryTimeDetectorGraphFlow.lean}.
The coefficient calculation is in \code{OrdinaryTimeDetectorZeroMode.lean}.
Neither assumes the trajectory-selector tests or a previously selected
invariant measure.

\subsection{The compact graph and its invariant probability}

Write $H$ for the uncentered primitive just constructed and
$z=\kappa H$. The compact graph is the image of the torus under
\[
 x\longmapsto
 \left(\frac{(\Gamma(x),1)}{\sqrt{1+|\Gamma(x)|^2}},
       \frac{(2z(x),1-|z(x)|^2)}{1+|z(x)|^2}\right).
\]
The first factor is positive central compactification; the second is
inverse stereographic projection. The image is nonempty and compact,
because the torus is nonempty and compact and both maps are continuous.

Compactness alone does not make this image invariant. The proof follows
an actual source orbit $a(s)=\Gamma(x+\omega s)$ and uses the reciprocal
central clock below. Here $2K$ is the clearing exponent; in the detector
specialization, $K$ is \code{sourceDegreeBound I}, not the primitive
budget $K_e$.
\[
 t(s)=\int_0^s (1+|a(u)|^2)^K\,du.
\]
Its derivative is at least one, so this clock is an increasing
diffeomorphism of the whole real line. In compact time, both $a$ and
$z$ acquire the same factor $(1+|a|^2)^{-K}$. Differentiating the central
and stereographic chart maps therefore gives the literal compiled
polynomial vector field. Uniqueness identifies this curve with the
canonical compact flow. Thus that flow preserves the graph image.

This reasoning does not assume that $\Gamma$ is injective, or that $H$
is a function of $\Gamma$. Both are functions of the original torus
parameter. Nor does it assert equivariance in unchanged time: the
reciprocal clock is part of the proof.

Restrict the compact flow to this invariant image and apply
Krylov--Bogolyubov. Pushing the resulting invariant probability into
the ambient compact source gives the graph full mass. Since the literal
reward is at least $5w_*/4$ at every graph point, its integral is at least
$5w_*/4>0$. In particular, the maximal invariant average $\beta$ is
positive for an unsolvable arithmetic instance.

The proof is divided into four reusable steps:
\begin{itemize}
\item \code{Dynamics/CompactInvariantGraph.lean}: the measure theorem;
\item \code{OrdinaryTimeCompactGraphTrajectory.lean}: chart differentiation
and flow uniqueness;
\item \code{OrdinaryTimePositiveGraph.lean}: reciprocal-clock invariance;
\item \code{OrdinaryTimeDetectorInvariantGraph.lean}: the actual detector.
\end{itemize}
The last step uses the raw least-denominator gain code and its proved
decoding identity.

\subsection{The negative branch uses the same field}

The module \code{OrdinaryTimeDetectorGap.lean} proves the negative branch:
$\beta\le-\varepsilon$ for a solvable arithmetic instance with the stated
nonzero-column condition. The argument works for every nonzero rational
gain and every rational margin. A hypothetical ergodic invariant
probability with positive finite-chart mass gives an ergodic conditional
base probability. If all its drift means vanished, ordinary-time
extraction would imply arithmetic unsolvability. Thus some conditional
drift mean is nonzero. The compact-accumulator obstruction then forces
an ergodic maximizing measure outside the finite chart, onto the central
boundary or accumulator infinity. There the literal reward is at most
$-\varepsilon$.

Specializing to the same raw gain code and the margin $w_*/4$ proves
the negative branch for the field used in the positive graph argument.
No identification with the earlier alternative gain is needed. Both
branches concern the actual compiled flow and reward, rather than an
abstract sign-gap interface supplied as a hypothesis.

\subsection{The literal gain feeds a computable homogeneous compiler}

The sign theorem and the code generator must describe the same field.
The module \code{OrdinaryTimeDetectorSign.lean} packages the two branches
into a semantic certificate and proves that its source field and reward
are literally the existing sparse-polynomial compiler applied to the
new raw gain and margin codes. The faithful homogeneous wrapper is then
GAS exactly when the arithmetic instance is solvable.

With the matrices and prime tuple fixed, replacing the integer
right-hand side is primitive recursive. Computing the least-admissible
gain denominator is total computable. Composing these operations with
the polynomial compiler and homogeneous wrapper gives a computable
field-code map. Its definition takes no proof of solvability, graph
invariance, or the sign theorem. The module
\code{OrdinaryTimeDetectorCompiler.lean} proves both its computability
and its literal equality to the field carried by the semantic certificate.

Finally, \code{OrdinaryTimeUniversalDetector.lean} applies this map to
the fixed augmented universal DQL template. Only the right-hand side
varies; the augmentation supplies the required nonzero linear columns.
All field dimensions and degree bounds are therefore fixed. This gives
an exact code-level equivalence between universal halting and HGAS for
the printed-gain family. The preserved RE upper bound and cubicization
and padding theorems give fixed-pair and cofinal
$\Sigma^0_1$-completeness. This composition does not identify the printed
gain with the earlier alternative gain; it verifies a separate effective
family using the same established polynomial compiler and wrapper.

\section{A fixed computational source and register instructions}

The module \code{Computability/UniversalMachineSource.lean} starts with
Mathlib's partial-recursive universal evaluator, not an assumed universal
register machine. A pair consisting of a program code and its input has
a primitive-recursive natural-number encoding. The list-code completeness
theorem supplies one fixed program computing that evaluator on singleton
lists. Mathlib's verified translation then gives a fixed stack-machine
interpreter. Its halting predicate agrees with \code{HGAS.UniversalHalts}
under the input encoding, so its halting language is $\Sigma^0_1$-complete.

The compiled interpreter has one finite control support independent of
the input. This support is derived from the compiler and contains every
reachable active label. The interpreter code is chosen from a proved
existence theorem; it is not an evaluated numeric program listing.

Separately, \code{Computability/RegisterMachine.lean} defines finite
increment/decrement programs with natural-valued registers, finitely many
active labels, and an absorbing halting label. Its total finite-step run
halts exactly when Mathlib's partial transition evaluator is defined.
The clear-register loop takes $n+1$ steps when the register initially holds
$n$: $n$ decrements and a final zero test. The transfer loop takes $2n+1$
steps: each decrement is followed by a target increment, then a final
zero test. The formal equations preserve all other registers.

The module \code{RegisterMachineArithmetic.lean} implements the arithmetic
needed for numeral stacks. Symbols use the nonzero digits $1,\ldots,b-1$,
with the top at the least significant digit; the empty stack has code
zero. Division by $b$ removes the top, remainder reads its digit, and the
encoding is injective. Four symbols can therefore use base five.

To push digit $d$, decrement the value register once and increment a
cleared scratch register $b$ times, repeating until the value is zero.
Transfer the scratch back, then increment $d$ times. This produces $bN+d$
and clears the scratch. To pop, cycle through $b$ decrement/zero-test
labels, incrementing the quotient scratch after each complete cycle.
The test at which zero is found records the remainder in finite control.
Its continuation transfers the quotient back and clears the scratch.
The correctness theorems derive these runs from literal instruction
equations and preserve every other register. They include the final zero
tests in their exact step counts; they do not assume the desired macro
behavior as part of an interface.

The modules \code{RegisterMachineStackBlocks.lean} and
\code{RegisterMachineDispatch.lean} now allocate finite push/pop blocks
and finite store-dispatch tables. Mathlib's interpreter uses four stacks,
four symbols, and a five-valued optional-symbol store. Four stack
registers and one cleared scratch register suffice for these operations.
The store selects an initial control label in a fixed program; it does
not select a different program at runtime. A push block has 12 internal
labels and a pop block 16; allocating all five store branches gives 60
and 80 labels, respectively. The proofs match the literal Mathlib
instruction semantics, including store updates and continuation labels.

\code{RegisterMachinePeek.lean} implements peek by pop followed by
restoration, retaining both the old store and popped symbol in finite
control. Its proof restores every stack and clears the scratch register,
including the empty-stack case. \code{RegisterMachineEmbedding.lean}
then allows a block's exit labels to refer to any ambient program labels,
including backward cyclic jumps. The proof reindexes an already verified
block whose continuation is the ambient program itself. Only literal
local instruction equations are required; a simulation property is not
assumed.

\code{RegisterMachineStatement.lean} recursively compiles all seven
constructors of Mathlib's actual stack-statement syntax. Both arms of a
branch receive fixed finite label ranges. A structural induction proves
that each complete statement runs for a positive finite number of
increment/decrement steps and reaches exactly the result of
\code{TM2.stepAux}, with its finite store encoded in the exit label and
the scratch register cleared. The external destination instruction has
not yet executed. This is an actual compiler and semantic theorem, not
an assumed interpreter interface.

\code{RegisterMachineFinite.lean} assembles these statements into one
finite cyclic program. The block width is computed from the maximum
statement size, so neither a size bound nor a simulation relation is
assumed. Its halting theorem compares the actual register run with
Mathlib's partial stack-machine evaluator in both directions.
\code{RegisterMachineSupport.lean} restricts the interpreter's control
labels to its proved finite support. Its arbitrary fallback outside that
support only makes the encoding total; reachable configurations never
use it.

Finally, \code{RegisterMachineInput.lean} encodes the actual binary input
stack. The primitive-recursive numeral map satisfies
\[
 F(0)=2,\qquad F(n)=5F(\lfloor n/2\rfloor)+3+(n\bmod2)\quad(n>0).
\]
All other stack registers and the scratch register initially vanish.
\code{RegisterMachineUniversal.lean} therefore obtains one fixed finite
five-register INC/DEC program, with the input in register zero, and an
explicit primitive-recursive reduction from \code{HGAS.UniversalHalts}
to its literal halting predicate. Universality is now a conclusion of
the assembled compiler.

The membership argument applies to every raw input numeral, not only
encoded source inputs. \code{RegisterMachineComputability.lean} proves
that a finite program's packed step and finite run are primitive recursive;
searching for a halting label semidecides its halting language. Together
with the reduction, this proves $\Sigma^0_1$-completeness of the one fixed
register program's input language. The later monitor compiler and
\code{UniversalStabilityReduction.lean} connect this language to HLS.

\section{Restart budgets}

The following finite inequality is the counting argument in
\cite[Proposition 4.5, Step 3]{korda}. It is the first completed component
of the new proof-based formalization of that reduction.

\begin{lemma}
Let $h_j$ be real budgets, let $e_j\ge0$ be errors, and distinguish some
turns as productive. Suppose, for $j<n$,
\[
 h_{j+1}\ge h_j-e_j,
 \qquad
 h_{j+1}\ge\max(h_j+1,1)-e_j
 \quad\text{when turn $j$ is productive}.
\]
If $P_n>0$ is the number of productive turns before $n$, then
\[
 h_n\ge P_n-\sum_{j<n}e_j.
\]
In particular, if $P_n\ge T\in\N$ and $\sum_{j<n}e_j<1$, then
$\max(h_n,0)+1>T$.
\end{lemma}

\begin{proof}
At the first productive turn the reset to at least one removes all
dependence on the initial budget, which may be negative. Each later
productive turn adds at least one before its error; every other turn loses
at most its error. Induction gives the displayed bound. If $P_n>0$, combine
it with $P_n\ge T$ and the total-error bound. If $P_n=0$, then $T=0$,
and $\max(h_n,0)+1\ge1>0$.
\end{proof}

The exact finite-prefix statement is
\code{Stability.Simulation.productiveCount_sub_error_le_budget}, in
\code{Stability/Simulation/RestartBudget.lean}. It does not assume a
simulation-correctness package. The clocked component estimates below
now supply the corresponding one-sided ODE inequalities; the full
monitor still needs their connection to its output events and discrete
halting runs.

\subsection{Literal discrete restart runs}

\code{RegisterRestart.lean} defines the restart machine itself, with a
natural budget, integer countdown, natural registers, and a restart,
active, or halt label. A restart increases the budget by one, reloads
the input, and begins a clean attempt. Its budget after any finite run
equals its initial budget plus the actual number of resets.

If the input machine halts within $T$ steps, an attempt funded with at
least $T$ units reaches halt before it can restart. Therefore every
finite run, from every state in this discrete domain, has at most $T$
reset labels: by the $T$th reset the newly granted budget is sufficient,
and no later reset is possible. This also covers arbitrary initial
registers and negative integer countdowns.

Conversely, if the input never halts, each clean funded attempt exhausts
its countdown and returns to restart. Starting at the designated zero
budget restart state gives arbitrarily many resets.
\code{RegisterRestartNonhalting.lean} proves this construction and the
equivalence between input halting and a uniform finite reset-count bound.
It does not claim unbounded counts from every initial state: a state
already at halt is an immediate counterexample to such a claim.

\code{RegisterRestartEmbedding.lean} proves that the explicit real
transition map intertwines with every discrete step and hence every
finite iterate. The reset-count equivalence therefore holds on the
embedded discrete domain. Extending the required tracking estimates to
arbitrary real states of the continuous monitor is a further task.

\section{A finite-interval cubic relaxation estimate}

The first analytic component of the simulation proof is a scalar
relaxation estimate \cite[Lemma 4.1]{korda}. The completed module
\code{Stability/Simulation/CubicRelaxation.lean} proves its strictly
positive-excess branch as a differential inequality, not as an assumed
simulation property.

\begin{lemma}
Let $a\le b$. Suppose $w$ and $\gamma$ are continuous on $[a,b]$, $w>0$
there, and $w$ is differentiable in $(a,b)$ with
\[
 w'(t)\le-\gamma(t)w(t)^3.
\]
Then
\[
 \frac1{w(a)^2}+2\int_a^b\gamma(t)\,dt\le\frac1{w(b)^2}.
\]
If $\int_a^b\gamma>0$, it follows that
\[
 w(b)<\left(2\int_a^b\gamma(t)\,dt\right)^{-1/2}.
\]
\end{lemma}

\begin{proof}
Since $w>0$, differentiation gives
$(w^{-2})'=-2w'/w^3\ge2\gamma$ in the interior. The integral comparison
theorem bounds the integral of the continuous lower derivative bound by
the endpoint difference of $w^{-2}$. This proves the first inequality
without requiring continuity of $w'$. The initial reciprocal is strictly
positive. Dropping it therefore gives a strict inequality, which yields
the second assertion after multiplying by $w(b)^2$ and taking square roots.
\end{proof}

The reciprocal inequality does not require $\gamma\ge0$; only the final
inverse-square-root estimate needs positive integrated gain. The important
restriction is $w>0$ on the whole interval. Handling zeros, moving targets,
additive forcing, and continuation requires further arguments. These
theorems are not labeled as a formalization of the entire relaxation lemma.

\subsection{A constant target, including sign changes and zeros}

The companion module
\code{Stability/Simulation/ConstantTargetCubicRelaxation.lean} handles an
actual unforced solution segment
\[
 y'=\gamma(t)(A-y)^3,\qquad \gamma\ge0,
\]
with constant target $A$, without any restriction on the initial value.
The squared error satisfies
\[
 \frac{d}{dt}(y-A)^2=-2\gamma(t)(y-A)^4\le0.
\]
Consequently, reaching the target at any time forces the solution to stay
there for the rest of the interval. If $y(b)\ne A$, the error cannot have
vanished earlier. Its absolute value is therefore differentiable throughout
the interior and satisfies
\[
 \frac{d}{dt}|y-A|=-\gamma(t)|y-A|^3.
\]
The positive-excess lemma applies. If $y(b)=A$, the same final bound is
immediate. Thus, for positive integrated gain,
\[
 |y(b)-A|<\left(2\int_a^b\gamma(t)\,dt\right)^{-1/2}.
\]
This proof uses no separate uniqueness theorem. It assumes a solution
segment is already defined. The next two modules extend the estimate to
moving targets and forcing; existence remains a separate obligation.

\subsection{Moving targets and additive forcing}

Suppose the supplied scalar solution satisfies
\[
 y'=\gamma(t)(A(t)-y)^3+\varphi(t),\qquad
 \gamma(t)\ge0,\qquad |A(t)-a_*|\le r,
\]
on a finite interval $[a,b]$, with continuous rate and forcing and with
positive integrated gain. Define the cumulative forcing mass
\[
 \Phi(t)=\int_a^t|\varphi(s)|\,ds.
\]
For the upper excess $w(t)=y(t)-(a_*+r)-\Phi(t)$, differentiation and
$\varphi-|\varphi|\le0$ give $w'\le-\gamma w^3$ whenever $w>0$.
The squared positive part $(\max(w,0))^2$ is differentiable even at
$w=0$, and its derivative is nonpositive. Thus an excess cannot cross
from nonpositive to positive. If its final value is positive, it was
positive throughout and the reciprocal-square estimate applies. If its
final value is nonpositive, the desired upper bound is immediate. Applying
the same argument to $(a_*-r)-y-\Phi$ yields
\[
 |y(b)-a_*|<r+\left(2\int_a^b\gamma(t)\,dt\right)^{-1/2}
            +\int_a^b|\varphi(t)|\,dt.
\]
There is only one forcing-mass term, not two.

The same no-crossing argument, without a positive-gain hypothesis, proves
\[
 |y(t)|\le\max\{|y(a)|,R\}+\int_a^b|\varphi(s)|\,ds
 \quad(a\le t\le b)
\]
whenever $|A(t)|\le R$. One-sided target bounds give the corresponding
one-sided relaxation estimates. These results are in
\code{MovingTargetCubicRelaxation.lean} and
\code{ForcedCubicRelaxation.lean} under \code{Stability/Simulation/}.
They formalize the finite-segment estimates in
\cite[Lemma 4.1]{korda}, but do not themselves construct a solution or
prove an ODE continuation theorem.

\section{Clock gain and a uniform summable error budget}

The module \code{Simulation/ClockGain.lean} starts with the actual clock
equations on a supplied finite interval:
\[
 \theta'=1,\qquad u'=-v,\qquad v'=u.
\]
The first equation gives $\theta(t)=\theta(a)+t-a$. Differentiating the
two rotating coordinates gives
\[
 \begin{split}
 u(t)&=u(a)\cos(t-a)-v(a)\sin(t-a),\\
 v(t)&=u(a)\sin(t-a)+v(a)\cos(t-a).
 \end{split}
\]
In particular, $u^2+v^2$ is constant. These are identities for supplied
solution segments, not assertions of global existence.

On a centered active arc, put
\[
 \gamma(t)=K(1+(c+t)^2)^2\max(r\cos t,0),
 \qquad |t|\le\pi/3,
\]
where $K>0$, $r\ge0$, and $r^2>1/2$. Set
$d=\max(|c|-\pi/3,0)$. Then $|c+t|\ge d$, while
$r>2/3$ and $\cos t\ge1/2$ give $\max(r\cos t,0)\ge1/3$.
Integrating this pointwise estimate, using $4\pi>9$, yields
\[
 K(1+d^2)^2\le2\int_{-\pi/3}^{\pi/3}\gamma(t)\,dt.
\]
Thus the relaxation error from the preceding section is at most
\[
 K^{-1/2}e(c),\qquad
 e(c)=\frac1{1+\max(|c|-\pi/3,0)^2}.
\]

The module \code{Simulation/ClockErrorSummability.lean} proves the uniform
bound
\[
 \sum_{j\in\Z}e(c+j\pi)\le\frac{29}{10}<3
 \qquad(c\in\R).
\]
To see the estimate, shift $c$ by an integer multiple of $\pi$ into
$[0,\pi]$. Split the integer lattice into $j=n$ and $j=-(n+1)$.
On either half, the distance of the $n$th point from zero is at least
$n\pi$, so its error is at most $a_n=(1+4n^2)^{-1}$. Here $a_0=1$,
$a_1=1/5$, and for $n\ge2$,
\[
 a_n\le\frac1{4n(n-1)}
       =\frac1{4(n-1)}-\frac1{4n}.
\]
The tail telescopes to at most $1/4$, giving
$\sum_{n\ge0}a_n\le29/20$ on each half. Summability is established
before manipulating infinite sums. The forward lattice and every finite
subset inherit the strict bound $3$, and multiplying by $K^{-1/2}$ gives
the error budget in \cite[Lemma 4.2]{korda}. Identifying these centers with
every turn of the actual monitor, including initial partial turns, requires
further bookkeeping.

The module \code{AlternatingClock.lean} supplies this connection for an
explicitly aligned phase $u(t_0)=0$, $v(t_0)=-r$. Its half-turn boundaries
are $\tau_n=t_0+n\pi$, and their centers are $c_n=\tau_n+\pi/2$.
The oscillator equation gives
\[
 (-1)^n u(c_n+s)=r\cos s,\qquad
 \theta(c_n)=\theta(t_0)+\pi/2+n\pi.
\]
The opposite buffer's gate vanishes throughout the corresponding
half-turn. The actual gain in centered coordinates is therefore the
gain function estimated above. Half-open half-turns are pairwise disjoint.

The module \code{HalfTurnForcingMass.lean} proves that the absolute
forcing integrals over the first $M$ half-turns sum exactly to the
integral over $[t_0,\tau_M]$. Combining this partition with the clock
series gives, in \code{HalfTurnErrorBudget.lean},
\[
 \sum_{n<M}\left(2\int_{\tau_n}^{\tau_{n+1}}|\varphi(t)|\,dt
       +K^{-1/2}e(\theta(c_n))\right)
 <2\int_{t_0}^{\tau_M}|\varphi(t)|\,dt+3K^{-1/2}.
\]
For $K=2^{20}$ and total forcing mass at most $1/256$, this is less than
$11/1024$.

\subsection{The actual compute/copy update}

Three modules connect the allowances to the actual two-buffer dynamics:
\begin{quote}
\code{DoubleBufferUpdate.lean}\\
\code{HalfTurnGain.lean}\\
\code{ClockedDoubleBufferUpdate.lean}.
\end{quote}
During a compute half-turn, the copy gate vanishes,
so the held buffer $a$ moves only by its forcing. Write $E_c$ for the
integral of the larger of the two forcing norms during this half-turn.
Then $\|a(t)-a(t_0)\|_\infty\le E_c$.

Suppose the real transition map agrees along this compute path with a
map $F$ of Lipschitz constant at most one. Its moving target stays within
$E_c$ of $F(a(t_0))$. The forced cubic-relaxation estimate gives
\[
 \|b(t_1)-F(a(t_0))\|_\infty\le2E_c+\rho_c,
 \qquad \rho_c=\left(2\int_{t_0}^{t_1}\gamma_c\right)^{-1/2}.
\]
During the following copy half-turn, $b$ is held and moves by at most
its forcing budget $E_p$. Applying the same estimate to $a$ gives
\[
 \|a(t_2)-F(a(t_0))\|_\infty
 \le2(E_c+E_p)+\rho_c+\rho_p.
\]
There is no hypothesis on the initial distance of $b$ from its target.

The whole-half-turn gain is at least the central-arc gain, by positivity
and interval inclusion. The clock identities therefore bound $\rho_c$
and $\rho_p$ by the corresponding lattice errors. The final theorem
starts with the oscillator-gated component ODEs and concludes that one
update error is at most the sum of its two half-turn allowances. It does
not assume vanishing gates, positive integrated gain, or the update
estimate.

\subsection{The real transition map, budgets, and clean resets}

The real transition map is now explicit in
\code{RegisterTransition.lean}. It uses a real budget, countdown, label,
and register coordinates. Integer active labels are encoded by
$1,\ldots,s$, restart by zero, and halt by $s+1$. Piecewise-linear gates
select the instruction branch and distinguish zero from positive
registers. On the stated regular neighborhoods of radius $1/8$, these
gates are constant. Freezing them gives an affine branch of Lipschitz
constant at most one in the sup norm. The reset branch has the same
bound because $h\mapsto\max(h,0)+1$ is nonexpansive; the halt branch
fixes honest halted states. These facts are proved in
\code{RegisterLocalBranches.lean}.

\code{RegisterTransitionExactness.lean} checks the discrete meaning:
with natural registers and countdown at least one, an active branch
executes the literal increment/decrement instruction. With countdown at
most zero it preserves the registers and selects restart. Restart sets
both budget and countdown to $\max(h,0)+1$, restores the input registers,
and selects the entry label. No claim about every fractional countdown
is hidden in these exactness statements.

\code{RegisterTransitionFormula.lean} expresses each coordinate of the
same map by an explicit absolute-value formula and proves equality of
their evaluations. Finite sums retain every branch, including zero
terms. Changing the input natural numbers therefore changes only
rational constant values, not the erased formula tree, radical count,
or majorant degree bound. This is the syntactic resource theorem used
later to prove primitive recursiveness of the raw-code compiler.

The modules \code{DoubleBufferBudget.lean} and \code{ClockedBudget.lean}
derive one-sided budget bounds directly from the component ODEs. The
ordinary bound uses the lower target $f(h)=h$; the reset bound uses
$f(h)=\max(h+1,1)$. Both functions have Lipschitz constant one, so the
same forcing and gain allowances apply.

Finally, the held buffer moves by at most one forcing mass between any
two times in a compute interval. If that mass is at most $1/16$ and the
label lies in $(-1/16,1/16)$ once, it stays in $(-1/8,1/8)$ throughout
the interval. The real map is then the reset map along the entire path,
and the compute/copy estimate proves a clean reset.
\code{HeldBufferMotion.lean} and \code{DoubleBufferReset.lean} contain
this argument.

\subsection{The constructed monitor and finite-phase bounds}

\code{ClockedMonitor.lean} assembles the actual field, with
$K=2^{20}$, and output
\[
 o=\max\{1-2|u^2+v^2-1|,0\}\,\max\{u,0\}\,\lambda(4a_{\rm label}),
\]
where $\lambda$ is the compactly supported label gate. Its positivity
implies $1/2<u^2+v^2<3/2$, $u>0$, and
$|a_{\rm label}|<1/16$. Thus the amplitude and reset-label restrictions
used above follow from an actual output event. The module
\code{ClockedMonitorFormula.lean} represents every field coordinate and
the output in the common absolute-value syntax, with input-independent
erased formula trees. \code{ClockedMonitorDynamics.lean} applies the
update and clean-reset estimates to the actual vector ODE with additive
buffer forcing.

\code{FundedRealRun.lean} handles real, not necessarily integer,
countdowns during a sufficiently funded clean attempt. If $B\ge T$, the
first $T$ iterates follow the actual register computation, with unused
countdown at least $B-T$. If the computation halts, the reached real
state is fixed. No claim is made that an exhausted fractional countdown
continues to represent an honest machine state.

\code{FundedMonitorTracking.lean} uses the finite compute/copy bounds to
track such a funded real attempt. The accumulated gain allowances and
actual forcing mass keep the sampled error below the regular-branch
radius. Thus each local branch used in the induction is justified by the
preceding error bound, rather than supplied as an extra hypothesis.
\code{HaltingMonitorTracking.lean} continues the comparison after the
source halts: the reference state is then fixed, so no further countdown
funding is needed. The held-buffer estimate also excludes positive output
throughout each post-halt turn, not only at sample times.

This yields an all-initial-buffer event count. Call a complete turn an
event if the actual output is positive somewhere during it. Each event
forces a clean reset. The ordinary budget recurrence prevents substantial
loss between events, while the productive recurrence increases the budget
at an event. If the source halts after $T$ instructions and the entire
segment's forcing mass is at most $1/256$, a $(T+1)$st event supplies a
sufficiently funded reset. Tracking from that reset excludes any later
event. Hence at most $T+1$ turns are events. No initial buffer is assumed
to encode a machine state. The phase is aligned, but the required
amplitude restriction follows from any alleged event itself.
\code{MonitorEvents.lean} and \code{MonitorEventBound.lean} prove this
statement for the actual forced ODE and actual output.

The output is nonnegative and at most $2$. Integrating over complete
turns therefore costs at most $4\pi$ per event. Aligning an arbitrary
phase discards at most two partial turns, giving the bound
$4\pi(T+3)$ whenever total forcing mass is at most $1/256$.
For arbitrary forcing, cut the interval whenever its accumulated mass
reaches $1/256$. Continuity of this actual cumulative integral supplies
the cuts by the intermediate value theorem. Adjacent-interval additivity
then gives
\[
 \int o\,dt\le 4\pi(T+3)(256M+2),
\]
where $M$ is the actual forcing mass.
\code{MonitorOutputIntegral.lean} and \code{MonitorForcingIntegral.lean}
prove these bounds without any assumption on initial buffers, phase,
amplitude, or segment length.

For nonhalting, start the held buffer at the designated integer restart
state and use unit oscillator amplitude. The other buffer may be
arbitrary. \code{IntegerMonitorTracking.lean} proves the real map's local
branches at every embedded integer restart state and tracks every finite
prefix under the small total forcing budget. On a reference reset turn,
the held label stays in the region where the output gate equals one.
The actual output is then the oscillator sine throughout the compute
half-turn, whose integral is $2$. Thus the integral is at least twice
the exact reference reset count. That count is unbounded for nonhalting.
The genuine global smoothed solution has forcing mass less than
$\epsilon\pi<1/256$ when $0<\epsilon\le1/1024$; hence this lower bound
applies to it. \code{NonhaltingMonitorOutput.lean} constructs the resulting
global witness rather than assuming its existence.

The unsmoothed field is locally Lipschitz, although absolute values need
not be differentiable. \code{FormulaLocalExistence.lean} proves this
regularity structurally and applies local ODE existence through every
state. \code{OscillatorPhase.lean} finds an aligned phase on any supplied
full period, including the zero-amplitude case. For shorter segments,
\code{MonitorPhaseBounds.lean} needs only a constant sign of $u$.
On a compute segment, $a$ moves by at most the forcing mass. The actual
continuous transition map $H$ is bounded on the resulting compact ball,
and cubic relaxation bounds $b$ by the larger of its initial norm and
that target bound, plus forcing mass. On a copy segment, the roles are
reversed. Hence a uniform bound for either kind of segment can be chosen
from the initial buffer bound and forcing-mass cap before seeing the
trajectory. \code{OscillatorPartition.lean} chooses a sine origin in the
period preceding the initial time and clips its half-turn grid to the
supplied interval. Every resulting segment has constant sign, even for
zero amplitude or a degenerate clipped segment. The number of pieces can
be fixed from a time-horizon cap. Iterating the uniform phase bound gives,
in \code{MonitorFiniteBounds.lean}, a full-state bound chosen from the
time horizon, initial norm, and forcing-mass cap before the trajectory.
No future solution is assumed to select this partition.

\code{BoundedContinuation.lean} turns these a priori estimates into
existence. For a prescribed horizon, multiply a locally Lipschitz field
by a Lipschitz norm cutoff which is one on a ball larger than the proved
bound and zero outside a slightly larger ball. The cutoff field is
bounded and globally Lipschitz, so it has a global trajectory. If that
trajectory first reached the inner boundary, its preceding segment would
solve the original equation and violate the a priori bound. It therefore
solves the original equation throughout the horizon. Local uniqueness
makes the trajectories constructed for different horizons agree, allowing
them to be joined into a global forward trajectory. The proof does not
require the original field to be differentiable.

\code{MonitorCompleteness.lean} applies this theorem to the actual
unsmoothed monitor, for every real initial state. Thus forward existence
is now proved, not inferred from a priori estimates alone.
\code{SmoothedMonitorCompleteness.lean} then proves forward completeness
for the actual shared-precision smoothing. Smoothing leaves the clock
and oscillator formulas unchanged. Its remaining residual is a continuous
buffer forcing whose joint mass is less than $\varepsilon\pi$ on every
finite solution segment. The same a priori bound and continuation theorem
therefore apply. Complete robust tracking remains separate.

\section{Smoothing formulas with polynomial error bounds}

An absolute-value formula is built from rational constants and coordinates
using addition, multiplication, and absolute value. For $L\ge1$, replace
each absolute-value node from the inside out by
$t\mapsto\sqrt{t^2+L^{-2}}$. Denote the original and smoothed values of a
formula $E$ by $E(x)$ and $\widehat E_L(x)$.

\begin{theorem}
The explicit structural recursion below constructs rational polynomials
$U_E,V_E$, independent of $L$, such that pointwise
\[
 U_E,V_E\ge0,\qquad |E|,|\widehat E_L|\le U_E,\qquad
 |\widehat E_L-E|\le V_E/L.
\]
\end{theorem}

\begin{proof}
A coordinate receives $(U,V)=(1+x_i^2,0)$; a rational constant $c$ receives
$(1+|c|,0)$. For a sum, add both bounds. For a product, use
\[
 U=U_1U_2,\qquad V=V_1U_2+U_1V_2,
\]
and the identity
\[
 \widehat E_1\widehat E_2-E_1E_2
 = (\widehat E_1-E_1)\widehat E_2
     +E_1(\widehat E_2-E_2).
\]
For an absolute-value node, use $(U,V)=(U_1+1,V_1+1)$.
Indeed,
\[
 |t|\le\sqrt{t^2+L^{-2}}\le|t|+L^{-1},\qquad L^{-1}\le1,
\]
and $\bigl||s|-|t|\bigr|\le|s-t|$. These inequalities establish each
inductive step, including nonnegativity of the two majorants.
\end{proof}

The theorem is \code{AbsFormula.majorants_correct} in
\code{Stability/Simulation/FormulaSmoothing.lean}; it follows
\cite[Lemma 5.1]{korda}. Since $L$ is universally quantified at each point,
the same statement applies when the precision depends on the state.

The module \code{FormulaResources.lean} bounds the degrees of both
majorants by a shape-dependent integer: coordinates contribute two,
constants zero, addition takes a maximum, multiplication takes a sum,
and absolute value leaves the bound unchanged. It also counts every
absolute-value node, without simplifying special constant values. Erasing
all constant values preserves both quantities, even when equal constants
at different positions are subsequently changed independently.

\subsection{One precision for the entire vector field}

For a finite family $E_i$, a clock coordinate $\theta$, and rational
$0<\varepsilon\le1$, define the rational polynomials
\[
 V=1+\sum_i V_{E_i},\qquad
 L=\varepsilon^{-1}(1+\theta^2)V.
\]
The module \code{FormulaPrecision.lean} proves $V\ge1$, $L\ge1$, and
\[
 |\widehat E_{i,L(x)}(x)-E_i(x)|
 \le\frac{V_{E_i}(x)}{L(x)}
 \le\frac{\varepsilon}{1+\theta(x)^2}.
\]
If $D$ is the largest of the shape-dependent majorant degree bounds,
then $\deg V\le D$ and $\deg L\le D+2$. Changing any rational constant,
including the precision parameter, does not change this upper bound.

The module \code{FormulaSmoothness.lean} proves global smoothness of the
smoothed family. At each absolute-value node, the radicand is
$t(x)^2+L(x)^{-2}>0$, because $L(x)\ne0$. Smoothness of square root away
from zero, composition, and structural induction finish the proof. The
reusable composition theorem permits any smooth nonvanishing precision,
not just the polynomial chosen above.

The actual vector difference has sup norm at most the same clock tail;
there is no dimension factor for this norm. Along a continuous supplied
curve whose clock has derivative one,
\[
 \int_a^b\frac{\varepsilon}{1+(c+t)^2}\,dt
 =\varepsilon\bigl(\arctan(c+b)-\arctan(c+a)\bigr)
 <\varepsilon\pi.
\]
Here $c$ is determined by the initial clock value and may be negative.
The module \code{SmoothingPerturbation.lean} applies this comparison to
the norm of the actual smoothing error on every finite forward segment.
It does not require an ODE in the other coordinates. A forward-complete
clock curve also has an integrable smoothing perturbation with total
norm integral at most $\varepsilon\pi$; the passage to the ray first
proves integrability from the uniform finite-interval bounds.

The polynomials and resource bounds are explicit recursive mathematical
objects. The raw-code modules prove primitive recursiveness for the fixed
formula templates and their input-dependent rational constants. The
following compiler identifies the radical count with the length of its
actual extracted program.

\section{Auxiliary radicals and the quadratic penalty}

The module \code{Simulation/PowerFraction.lean} handles the local algebra
of radical extraction. Represent each rational expression by a polynomial
numerator $p$ and a denominator power $L^k$. Addition uses the common
denominator $L^{k+\ell}$, and multiplication adds the exponents; neither
operation cancels factors based on coefficient values. For a smoothing
node, set
\[
 D=p^2L^2+L^{2k}.
\]
If $L\ge1$, then $D\ge1$, including the case $k=0$. For $L>0$,
\[
 \sqrt{(p/L^k)^2+L^{-2}}=\frac{\sqrt D}{L^{k+1}}.
\]
This exact identity introduces no further approximation. The positivity
of $L$ is necessary to choose the displayed square-root branch.

\subsection{An ordered compiler and its finite coordinates}

\code{RadicalProgram.lean} represents a program by its ordered list of
polynomial radicands. A base variable is tagged separately from a root
index. Triangularity means that the radicand producing root $j$ uses only
roots with indices below $j$. Starting with arbitrary root values, the
evaluator successively overwrites root $j$ by the nonnegative square root
of its radicand. Later writes cannot change an earlier radicand's value.
Induction therefore proves every root equation in the \emph{final}
assignment, not merely at the moment that root was computed.

\code{FormulaRadicals.lean} compiles a formula from the inside out.
Addition and multiplication concatenate their two child programs;
an absolute-value node appends the radicand $D$ above and returns its
fresh root divided by $L^{k+1}$. The correctness theorem proves that
actual sequential evaluation followed by evaluation of the returned
fraction equals the smoothed formula. No root equations are assumed in
that theorem. The number of roots and the denominator exponent both
equal the formula's syntactic absolute-value count, even when special
coefficient values would permit cancellation.

\code{FormulaRadicalFamily.lean} applies the same construction to a
finite list of formulas, allocating consecutive disjoint root ranges.
One execution then represents all vector coordinates and any output
formula included in the list. Every emitted radicand is at least one
at \emph{every} extended real state when $L\ge1$, not only on the
canonical root graph.

For completeness, \code{FormulaRadicalResources.lean} gives an explicit
common degree bound. If $d$ bounds the degree of $L$ and $c(E)$ counts
absolute-value nodes, define
\begin{align*}
 B(x_i)&=1, & B(a)&=0,\\
 B(E+F)&=B(E)+B(F)+d(c(E)+c(F)),
 &B(EF)&=B(E)+B(F),\\
 B(|E|)&=1+2B(E)+2d(c(E)+1).
\end{align*}
Induction bounds the output numerator and every emitted radicand by
$B(E)$. Taking the maximum gives a family bound. Erasing or changing
rational constant values leaves this bound unchanged.

Natural indices simplify concatenation, but the final polynomial system
needs finitely many coordinates. \code{RadicalFiniteCoordinates.lean}
explicitly substitutes the corresponding finite variable for each index
below the root count and zero for any index outside it. The dependency
bound proves that no out-of-range variable occurs. Renaming back therefore
recovers each polynomial exactly, preserving its evaluation and total
degree. This construction also handles a formula family with no radicals.

\code{FiniteRadicalFamily.lean} packages the resulting finite polynomial
tuple and defines its canonical root graph by the actual sequential
evaluator. Its equations, output values, and degree bounds follow from
the preceding compiler proofs. \code{RadicalSmoothness.lean} proves
smoothness of this graph by induction over the root program: each
coordinate is a positive square root of a polynomial in already smooth
coordinates. Polynomial evaluation uses only finitely many coordinates,
so the intermediate natural-index representation causes no analytic
infinite-dimensional assumption. Taking one more positive square root
gives smooth fourth-root coordinates $w_j$ with $w_j^2=\xi_j$.

\code{RadicalLiftPolynomials.lean} makes the substitution
$\xi_j\mapsto w_j^2$ in the actual rational polynomials. The resulting
radicands $\mathcal D_j(x,w)$ are globally at least one, and the
constructed canonical graph satisfies $w_j^4=\mathcal D_j(x,w)$.
\code{PolynomialCurveDerivative.lean} proves the polynomial chain rule
on arbitrary finite coordinate types, so tagged base and root coordinates
need no chosen enumeration.

\subsection{Constructing the rational vector field}

\code{PolynomialFraction.lean} keeps each fraction as an explicit pair
of rational polynomials. Addition uses the product denominator, without
cancellation; multiplication and finite sums are equally explicit.
This makes every future denominator-clearing factor visible.
\code{RadicalDependencies.lean} proves that $\mathcal D_j$ depends only
on base variables and auxiliaries with indices less than $j$. In
particular, its current and later formal partial derivatives vanish.

Let $F$ be the constructed base fractions, and suppose velocities for
the earlier auxiliaries have been constructed. Define
\[
 V_j=\frac{w_j}{4\mathcal D_j}
 \left(\sum_i\partial_{x_i}\mathcal D_j\,F_i
       +\sum_{h<j}\partial_{w_h}\mathcal D_j\,V_h\right).
\]
\code{TriangularLift.lean} performs this finite recursion on polynomial
fractions. Earlier coordinates are unchanged by later stages. Therefore
the displayed numerator is the full directional derivative of
$\mathcal D_j$ along the final vector field: the omitted partials are
zero. The global bounds $L\ge1$ and $\mathcal D_j\ge1$ imply that every
constructed denominator is at least one, including far from the graph.
The resulting rational field is consequently smooth everywhere.

\code{CompiledFormulaLift.lean} instantiates these facts with the
compiler's actual radicands and output fractions. Each base coordinate
selects its original formula-list position. On the generated graph its
velocity is exactly that formula's smoothed value. No triangularity,
denominator, or simulation assertion about the compiled system is left
as an independent hypothesis; only the precision lower bound is needed.

\code{TriangularLiftInvariant.lean} now applies the full polynomial
chain rule to an actual solution segment. It proves conservation of
$w_j^4/\mathcal D_j$ and hence preservation of every exact quartic
constraint. This is stronger than merely assuming a differential
equation for each radicand in a scalar lemma.

Constraint preservation alone would not construct a lifted trajectory.
\code{CompiledGraphTrajectory.lean} supplies that missing direction.
Given a trajectory $x(t)$ of the smoothed base field, set
$y(t)=(x(t),w_*(x(t)))$, where $w_*$ is the actual positive root evaluator.
This curve is differentiable. Differentiating its exact identities
$w_j^4=\mathcal D_j$ determines its auxiliary velocities in order.
Indeed, after the earlier velocities are identified, the differentiated
equation is
\[
 4w_j^3\dot w_j=\dot{\mathcal D}_j.
\]
Multiplying by $w_j$ and using $w_j^4=\mathcal D_j>0$ gives exactly the
constructed fraction $V_j$. This triangular tangent-uniqueness argument
proves that $y$ solves the actual rational lifted ODE. It does not assume
a lifted solution or claim global existence of the base solution.

\code{ClockedMonitorLift.lean} makes this construction for the monitor's
actual field and output formulas, with their shared precision polynomial.
One fixed finite ordering allocates all roots. The output fraction on the
graph equals the actual smoothed output, and a smoothed monitor trajectory
lifts to the actual rational field. \code{ClockedMonitorResources.lean}
proves that the exact root count and the precision and finite-radicand
degree bounds do not change with the loaded input or precision coefficient.
The final penalized homogeneous degree bound also uses the resource
estimates below.
The resource modules now also follow the actual recursive rational
velocities. If all currently stored numerator and denominator degrees are
at most $b$, the coordinate count is $N$, and the radicand degrees are at
most $d$, one root update has the conservative degree bound
$(N+1)(d+b)+1$. Iterating this bound for the fixed number of roots controls
the actual stored fractions, without cancellation or dependence on the
loaded input. \code{ClockedMonitorClearing.lean} applies the resulting
bounds to joint field/output denominator clearing.

\code{PositiveDenominatorClearing.lean} performs the subsequent algebraic
operation explicitly. For fractions $p_i/q_i$, put
\[
 Q=\prod_i q_i,\qquad P_i=p_i\prod_{j\ne i}q_j.
\]
Then $Q\ge1$ when each $q_i\ge1$, and $P_i=Q(p_i/q_i)$ at every state.
The degree bounds are the corresponding sums of numerator and denominator
degrees. Field coordinates and output can share this multiplier. The
module proves these polynomial identities. The finite-curve time-change
and output-integral equivalence are proved separately below.

\code{LiftedMonitorExistence.lean} combines smoothed forward completeness
with graph transfer. An affine test curve identifies the derivative of
the graph map on the smoothed base velocity; the chain rule then gives
the transfer on arbitrary time domains, including the one-sided derivative
at zero. Consequently every canonical graph point starts an actual global
forward trajectory of the rational lift. Neither off-graph completeness
nor completeness after polynomial time acceleration is asserted.

\subsection{Invariants and control of off-graph errors}

The scalar identity behind the radical lift is proved in
\code{Simulation/RadicalInvariant.lean}. Suppose $D>0$ and a supplied
solution segment satisfies
\[
 D'=\Lambda,\qquad w'=\frac{w\Lambda}{4D}.
\]
Differentiating gives
\[
 \left(\frac{w^4}{D}\right)'
 =\frac{4w^3w'D-w^4D'}{D^2}=0.
\]
Consequently $w^4/D$ is constant. An exact initial constraint $w^4=D$
persists, and $w^2=\sqrt D$ because $w^2\ge0$. This avoids a spurious
negative square-root branch without imposing a sign condition on $w$.
The proof uses no division by $w$, so it also covers $w=0$ for inexact
initial constraints. It does not assert that those inexact constraints
decay.

When $D\ge1$, factorization also gives the global estimate
\[
 |w^2-\sqrt D|\le|w^4-D|,
\]
because the other factor $w^2+\sqrt D$ is at least one. For $D,E\ge1$,
the identity $D-E=(\sqrt D-\sqrt E)(\sqrt D+\sqrt E)$ gives
$2|\sqrt D-\sqrt E|\le|D-E|$. These are the elementary defect
comparisons used by the following triangular construction.

\subsection{Constructing the global defect weight}

Write $\rho_j(x)$ for the canonical positive root of the original
radicand $D_j(x,\rho_{<j})$. These are the square roots represented by
$w_j^2$, not the fourth-root coordinates $w_j$ themselves. We first need
polynomial bounds for the canonical roots. For a rational polynomial
$p$, let $m(p)$ be the sum of the absolute values of its coefficients.
On a coordinate cube of radius $R\ge1$,
\[
 |p(z)|\le m(p)R^{\deg p}.
\]
Suppose polynomial majorants $S_i(x)\ge|\rho_i(x)|$ have already been
constructed for $i<j$. Use
\[
 R_j(x)=1+\sum_a x_a^2+\sum_{i<j}S_i(x)^2,
 \qquad S_j(x)=1+m(D_j)R_j(x)^{\deg D_j}.
\]
The radius dominates every base coordinate and every earlier root.
Later roots can be set to zero because $D_j$ does not depend on them.
The inequality $\sqrt t\le1+|t|$ proves the next majorant.
\code{RadicalMajorant.lean} implements this finite recursion literally,
including the empty system.

Now fix any auxiliary assignment $w$, however far from the graph.
Both $w_i^2$ and $\rho_i(x)$ lie in the interval bounded by
$w_i^2+S_i(x)$. A single polynomial radius $R(x,w)\ge1$ therefore contains
both extended points $(x,w^2)$ and $(x,\rho(x))$. Applying the scalar
mean-value theorem along their joining segment gives
\[
 |p(x,w^2)-p(x,\rho)|
 \le\sum_i C_{p,i}(x,w)|w_i^2-\rho_i|,
 \qquad C_{p,i}=m(\partial_{\rho_i}p)R^{\deg p}.
\]
These weights are nonnegative rational polynomials. If a root does not
occur in $p$, its formal partial derivative is zero and its weight is
identically zero. The formal proof uses coordinatewise partials rather
than an unspecified Lipschitz constant; this retains triangularity.
The modules are \code{PolynomialDifferenceBounds.lean} and
\code{RadicalComparison.lean}.

Define the actual quartic defects and their norm by
\[
 r_j=w_j^4-D_j(x,w^2),\qquad E=\left(\sum_jr_j^2\right)^{1/2}.
\]
The two elementary square-root inequalities above and the polynomial
comparison yield, with $d_j=|w_j^2-\rho_j(x)|$,
\[
 d_j\le |r_j|+\sum_{i<j}C_{D_j,i}d_i.
\]
We have discarded a factor $1/2$ to simplify the bound. Define actual
polynomials recursively by
\[
 W_j=1+\sum_{i<j}C_{D_j,i}W_i.
\]
Since $|r_j|\le E$, induction gives $d_j\le W_jE$. Thus any numerator
has error at most $(\sum_iC_{p,i}W_i)E$.
\code{TriangularErrorWeights.lean} isolates the reusable substitution
argument, and \code{RadicalDefectBounds.lean} supplies its actual
coefficients and defects. There is no smallness assumption on $E$.

Each compiled output divides its numerator by a common precision power
$L^e\ge1$. This cannot enlarge the error. Sum the finitely many output
weights and add $1$ to obtain one polynomial $W\ge1$ controlling every
field coordinate and the output. \code{CompiledFormulaError.lean}
instantiates all root equations and positivity hypotheses with the
compiler's proved properties. The associated resource modules bound
these actual weights' degrees using only the original degree bounds and
root count, not the coefficients.

\code{CompiledPolynomialFormula.lean} proves a useful exactness fact:
an abs-free formula allocates no roots, so its extracted value is correct
for every auxiliary assignment, even inside a family containing other
formulas with radicals. Therefore the lifted monitor still satisfies
$\theta'=1$, $u'=-v$, and $v'=u$ off graph. The buffer and output errors
are bounded by the constructed $WE$.

Finally, \code{ClockedMonitorPenalty.lean} constructs the polynomial
\[
 P=(1+\theta^2)W^2\sum_jr_j^2=(1+\theta^2)(WE)^2
\]
and subtracts it from the actual rational output, keeping its denominator
unchanged. The penalty is nonnegative everywhere and zero on the
constructed graph. Consequently the penalized output there is exactly
the smoothed monitor output. This is a pointwise algebraic theorem;
the following integral argument supplies its dynamical meaning.

The module \code{Simulation/QuadraticPenalty.lean} proves the separate
numerical estimate which later absorbs the weighted defect cost. For
every $a>0$ and real $c,R$, completing the square gives
\[
 cR-aR^2\le\frac{c^2}{4a}.
\]
Along a unit-speed clock, use $a=1+\theta^2$ and integrate. For every
continuous error function $R$ on a supplied finite forward interval,
\[
 \int\bigl(cR-(1+\theta^2)R^2\bigr)\,dt
 \le\frac{\pi c^2}{4}.
\]
The bound does not depend on the size of $R$ or the initial clock value.

\subsection{The rational-monitor criterion}

Set $R=WE$ and $\sigma=\epsilon/(1+\theta^2)$. Along any finite trajectory
of the actual rational lift, its base projection satisfies the original
monitor equation with a continuous buffer forcing of norm at most
$\sigma+R$. This is an identity obtained by subtracting the two actual
fields; the clock and oscillator components cancel exactly. The rational
output differs from the original monitor output by at most $\sigma+R$.
Consequently, if the source halts in $T$ instructions and
$C=4\pi(T+3)$, the proved monitor estimate gives
\[
 \int o\,dt\le C\left(256\left(\epsilon\pi+\int R\,dt\right)+2\right).
\]
Subtract the actual penalty from the rational output. With $c=256C+1$,
the preceding square-completion inequality yields the uniform bound
\[
 \int g\,dt\le 2C+c\epsilon\pi+\frac{\pi c^2}{4}.
\]
It holds for every finite solution segment, regardless of initial state
or length. The proof requires neither a small defect nor global
existence away from the graph. The residual identity is in
\code{ClockedMonitorLiftForcing.lean}. Its mass bound is proved in
\code{ClockedMonitorLiftIntegral.lean}, and the square-completion argument
is in \code{PenalizedMonitorIntegral.lean}.

Conversely, the actual global smoothed nonhalting witness lifts to the
canonical rational graph. Its penalty is zero, and the total difference
between original and smoothed output on any finite interval is at most
$\epsilon\pi$. Its finite penalized-output integrals are therefore
unbounded. \code{NonhaltingLiftedOutput.lean} proves this statement for
the constructed trajectory.

Let $\mathcal B(F,g)$ mean that one real constant bounds
$\int_0^T g(y(t))\,dt$ for every $T\ge0$ and every continuous solution
segment of $y'=F(y)$ on $[0,T]$. The endpoint
\code{RationalMonitorCriterion.halts_iff_bounded_penalizedOutput} proves
\[
 \text{the source halts}\quad\Longleftrightarrow\quad
 \mathcal B(\text{actual rational field},\text{actual penalized output})
\]
for $0<\epsilon\le1/1024$. All field, output, and simulation data are
constructed, not supplied as hypotheses.

\subsection{Clearing denominators on finite trajectories}

The common denominator $Q$ is a polynomial with rational coefficients
and $Q\ge1$ on the whole state space. Set
\[
 P=QF,\qquad b=Qg,
\]
where $F$ is the rational field and $g$ its penalized output. Both $P$
and $b$ are polynomials. The penalty did not change the stored output
denominator, so the same $Q$ clears the field and the output.

Here it is important to compare finite trajectories, not to assume that
the polynomial field is complete. Given a trajectory $y$ of $F$ on
$[0,T]$, define
\[
 s(t)=\int_0^t\frac{du}{Q(y(u))}.
\]
Continuity and positivity make this clock strictly increasing, with a
continuously differentiable inverse on the interior of its finite image
$[0,S]$. The reparametrized curve $z(s)=y(t(s))$ satisfies $z'=P(z)$,
and change of variables gives
\[
 \int_0^S b(z(s))\,ds=\int_0^T g(y(t))\,dt.
\]
Starting with a finite trajectory of $P$ gives the reverse construction.
Thus exactly the same constant bounds the output integrals for either
system. Neither direction requires a bound on $Q$ outside the compact
image of the supplied trajectory.

The reusable result
\code{hasBoundedFiniteOutput_iff_timeScaled} proves this equivalence for
continuous positive speeds on a real normed space. Its finite-clock
construction extends the speed by positive constant tails, obtains a
global inverse clock, and then restricts to the required finite interval.
This device handles the endpoints without extending the trajectory as
an ODE solution. Zero-length intervals are included.

Consequently,
\code{ClockedMonitor.halts_iff_bounded_polynomialOutput} proves
\[
 \text{the source halts}\quad\Longleftrightarrow\quad\mathcal B(P,b)
\]
for the actual constructed polynomial pair. Its degree bounds depend
only on the fixed register program, not on the loaded input.

\subsection{From bounded output to homogeneous stability}

Let $P$ be a polynomial vector field and $b$ a polynomial scalar output
on $\R^n$. Choose an odd integer $\ell\ge1$ such that every term of $P$
has degree less than $\ell$ and every term of $b$ has degree less than
$\ell-1$. For $Y=(y,z)\in\R^{n+1}$, homogenization defines polynomials
\[
 A(Y)=(z^\ell P(y/z),0),\qquad
 \beta(Y)=z^{\ell-1}b(y/z).
\]
These expressions initially describe the finite chart $z\ne0$; the
polynomial definitions make sense at $z=0$ as well. Set
\[
 F(Y)=|Y|^2A(Y)-(Y\cdot A(Y))Y+|Y|^2\beta(Y)Y.
\]
This is homogeneous of odd degree $\ell+2$. The strict degree bounds
imply $A=0$ and $\beta=0$ on the equator $z=0$.

On the Euclidean unit sphere, the angular vector field and radial
reward of $F$ are
\[
 A(Y)-(Y\cdot A(Y))Y,\qquad \beta(Y),
\]
respectively. The preserved angular stability theorem says that the
origin is Lyapunov stable, with existence, exactly when one constant
bounds all forward integrals of this reward along sphere trajectories.
The coordinate-space formulation uses the algebraic sum of squares for
the sphere; it does not replace that sphere by the sup-norm unit sphere.

Every equator point is stationary and contributes zero reward. An angular
trajectory starting outside the equator cannot reach it in finite time:
backward uniqueness would identify the trajectory with the stationary
equator solution. Thus the quotient $x=y/z$ is defined on every finite
segment of such a trajectory. Differentiation gives
\[
 x'=z^{\ell-1}P(x),\qquad \beta(Y)=z^{\ell-1}b(x).
\]
Since $\ell-1$ is even, the common speed is positive on either hemisphere.
Finite time change therefore transfers any uniform source-output bound
to the angular reward integrals.

For the converse, start with any finite trajectory of $x'=P(x)$ and
form its positive central lift
\[
 c(x)=(1+|x|^2)^{-1/2},\qquad Y=(c(x)x,c(x)).
\]
The lift lies on the unit sphere. Its velocity is the angular field
multiplied by $c^{-(\ell-1)}$, and its radial reward multiplied by the
same factor is exactly $b(x)$. Removing this positive speed gives a
finite angular trajectory with the same output integral. Finite-interval
uniqueness identifies it with the canonical angular flow used in the
stability criterion. This proves
\[
 \text{$0$ is Lyapunov stable for $F$}
 \quad\Longleftrightarrow\quad \mathcal B(P,b).
\]
Neither implication assumes global completeness of $P$.
The module and theorem are
\begin{center}
\code{ProjectiveStabilityCriterion.lean},\\
\code{lyapunovStable_fieldCode_iff_boundedFiniteOutput}.
\end{center}

The projective field itself is compiled by primitive-recursive sparse-code
operations. The stated degree conditions concern raw retained terms,
not just decoded degree: canceled high-degree terms must first be removed.
\code{PolynomialCodeTrim.lean} proves an effective filter that preserves
the decoded polynomial and converts its semantic degree bound into the
required raw-term bound. If both decoded degrees are at most $d$, the
module \code{PolynomialOutputReduction.lean} chooses $\ell=2d+3$ after
trimming. It produces a homogeneous field of degree $2d+5$ in $n+1$
variables, primitive-recursively from the original $n$-variable field and
output. The strict inequalities at the equator follow from this choice.
\subsection{Compiling the actual register monitor}

The module \code{ClockedMonitorPolynomialCode.lean} now constructs that
field/output pair from the loaded register values. Its correctness proof
identifies both decoded polynomials with the semantic monitor, not merely
with polynomials having the same degree bounds. The compiler first
instantiates fixed formula trees, extracts the ordered radicals, and
forms the literal rational fractions for the base and auxiliary
velocities. It computes the actual polynomial comparison weights and
quartic-defect penalty, clears the field and output denominators together,
and subtracts the same clearing factor times the penalty from the output.
Each operation is primitive recursive on sparse rational codes.

Only the rational coefficients depend on the loaded input. The coordinate
enumeration, number of roots, formula order, and degree bound $d$ are fixed
by the program. Equality casts identify the input-dependent semantic root
type with this fixed code space; they do not choose a new runtime
enumeration. Coefficient mass and degree are computed after coefficient
cancellation, so redundant sparse presentations do not invalidate the
bounds or the strict inequalities used in homogenization.

Fix the smoothing accuracy to $\varepsilon=1/1024$. Composing the exact
decoder, the polynomial-monitor criterion, finite-coordinate conjugacy,
and the projective theorem gives
\[
 \text{the fixed program halts on the loaded input}
 \quad\Longleftrightarrow\quad
 \mathrm{HLS}_{\mathbb Q}(n+1,2d+5)(F_{\mathrm{input}}).
\]
The module \code{HomogeneousStabilityReduction.lean} proves this equivalence,
homogeneity of every output code, and primitive recursiveness of the map.
The theorem has no supplied simulation, source-hardness, or effectivity
hypothesis. Its stability notion is exactly the preserved
\code{HGAS.HLSQ}, including forward existence near the origin.

\subsection{Unconditional cubic and cofinal HLS hardness}

Apply this compiler to the single five-register program derived earlier
from Mathlib's universal evaluator. The initialization puts the input in
register zero and sets the other registers to zero. This yields fixed
integers $n_0\ge2$ and $k_0\ge1$, and a primitive-recursive raw polynomial
map $f$, such that
\[
 \mathrm{UniversalHalts}(c,a)
 \quad\Longleftrightarrow\quad
 \mathrm{HLS}_{\mathbb Q}(n_0,2k_0+1)(f(c,a)).
\]
Both integers are closed symbolic terms independent of $(c,a)$; the
development does not evaluate them to numerical literals.

The preserved cubicization theorem is a computable many-one reduction
from this HLS language to degree three in the fixed dimension $N_0$,
given by \code{cubicLiftCodeDimension} at $(n_0,k_0)$. Its source-dimension
hypothesis is discharged because the monitor has a clock coordinate and
homogenization adds another coordinate. Dimension padding and odd-degree
raising then give every $N\ge N_0$ and every odd $d\ge3$.

The final module, \code{UniversalStabilityReduction.lean}, proves that
each resulting language is many-one hard for recursively enumerable
sets, is not computable, and has a complement that is not recursively
enumerable. The formal statements have no source-hardness assumption.
They do not assert that HLS itself is recursively enumerable, nor do
they settle its low-dimensional cases. The earlier conditional transfer
theorems remain unchanged and are used in this composition.

\section{Validation and limits}

The source-level ordinary-time equivalence, explicit Fourier and primitive
bounds, printed-gain compact graph measure, and two-sided reward gap are
proved. The halting-to-HLS reduction is also complete, from Mathlib's
universal evaluator through the actual rational and polynomial monitor
codes to the homogeneous field. Its final hardness theorems have no
source-hardness, simulation, or compiler-effectivity premise.

Independent audits checked the statements and the main semantic
connections: finite-time output estimates, arbitrary off-graph states,
positive denominator clearing, finite-curve projective charts, code
realization, graph invariance under the reciprocal clock, and both sides
of the compact reward gap. Regression examples include a one-step halting
program, an infinite increment loop, canceled polynomial terms,
zero-dimensional Fourier and chart cases, and a noninjective compact
graph parametrization.

The command \texttt{lake build Stability} builds the additions.
Two scripts provide separate checks:
\begin{itemize}
\item \code{scripts/StabilityRegression.lean} checks the examples;
\item \code{scripts/StabilityAxiomAudit.lean} reports axiom dependencies.
\end{itemize}
The policy remains: no project axioms, placeholders,
or unsafe proof shortcuts. Standard dependencies on \code{propext},
\code{Classical.choice}, and \code{Quot.sound} are permitted and reported.
The published build and import-closure audit remain separate and continue
to check the unchanged 837-module source tree. The ledger records the
validation commands and current counts.

These results do not determine the three-dimensional cubic problems,
evaluate the symbolic fixed dimensions, or prove recursive enumerability
of HLS or general polynomial GAS/LAS. They also do not claim that every
printed intermediate statement is exported verbatim: the ledger
distinguishes exact statements from equivalent formulations and lists
results outside this formalization's scope.

\begin{thebibliography}{9}
\bibitem{paper}
I.~O.~Shevchenko, J.~Liu, and X.~Liu,
\emph{Global asymptotic stability of homogeneous polynomial vector fields
is undecidable}, arXiv:2610.12434.
\url{https://arxiv.org/abs/2610.12434}.
For the intermediate statements discussed here, the precise working
manuscript is Overleaf commit \code{6f5bced}, especially
\code{sec08_fourier_to_orbits.tex}, \code{sec09_orbits_to_average.tex}, and
\code{sec11_proof_of_main.tex}.

\bibitem{korda}
M.~Korda,
\emph{Lyapunov stability of polynomial vector fields is undecidable},
arXiv:2609.22058v2.
\url{https://arxiv.org/html/2609.22058v2}.
The relevant statements are Definition 1.1 and Lemma 1.5 (stability with
existence and homogeneous linear bounds), Sections 3--4 (restarts and the
robust monitor), Lemma 4.1 (cubic relaxation), Proposition 4.5 (the
all-initial-state monitor estimate), and Sections 5--6 (algebraization and
homogeneous lifting).
\end{thebibliography}
\end{document}
